% =============================================================================
% Cell model: overview, conventions and Hill-gate summary.
% The per-module subsections (m0 ... m11, mmem) are \input at the end.
% Parameter values are those of dicty_reactions.jl at repository revision 2668d66
% (default configuration, all environment overrides unset).
% =============================================================================
\section{Intracellular Signalling Model}\label{sec:model}

This section specifies the reaction network of a single \emph{Dictyostelium discoideum} cell that is embedded in the reaction--diffusion simulation of the cell population.
For every module we list the reactions and constants in tables, explain why each reaction is present and which experiments motivate it, justify the constants, and justify every Hill gate.
The numbering of the modules is that of the simulation code; Modules~5--7 (Rac/PAK, PLA2/Ca$^{2+}$ and Ca$^{2+}$ influx) are specified in the model description but are not part of the network analysed here.

\subsection{Overview and conventions}\label{sec:conv}

\paragraph{Signal flow.}
Extracellular cAMP is sensed by the receptor cAR1 (Module~1), which catalyses the G-protein cycle (Module~2).
Free G$\beta\gamma$ drives two Ras branches: RasG, whose adaptation and amplification (Module~3) control PI3K and PIP3 (Module~4), and RasC, which acts through TORC2, the kinases PKB and the effector CRAC (Module~8).
Module~8 injects an activator into an excitable element (Module~0) that opens the adenylyl cyclase ACA (Modules~0 and~9) and thereby produces the cAMP pulse.
Intracellular cAMP is exported (Module~9), degraded intracellularly and feeds back through PKA, ERK2 and RegA (Module~10), and is degraded extracellularly by the secreted phosphodiesterase PdsA and its inhibitor PdiA (Module~11), which closes the loop.

\paragraph{Spatial discretisation and stochastic engine.}
The domain is a square lattice with voxel edge $2\,\mu$m, so that one voxel has a volume of $8\,\mu\mathrm{m}^3$ and $1\,\mu\mathrm{M}$ corresponds to $M=4817.6$ molecules per voxel.
A cell occupies $N_\mathrm{vox}=13$ voxels (a $3\times3$ block plus one voxel on each side), i.e.\ $104\,\mu\mathrm{m}^3$.
In every time step of $\Delta t=2$\,ms, each voxel runs an exact stochastic simulation of its reactions up to $\Delta t$~\cite{gillespie1977}, after which molecules move to neighbouring voxels; the diffusion step requires $D\Delta t/\Delta x^2\le0.25$, which holds for the fastest species, cAMP ($D=350\,\mu\mathrm{m^2/s}$, value $0.175$).
Membrane-bound and cytosolic species of the cell are confined to its voxels, whereas extracellular species (cAMP$_e$, PdsA$_e$, PdiA and their complex) live on the world lattice, which they leave at the domain boundary.
Because the cost of a step scales with the total propensity, uniform rate multipliers change the run time.

\paragraph{Units and rate constants.}
First-order constants are given in s$^{-1}$, bimolecular constants in $\mu\mathrm{M}^{-1}\mathrm{s}^{-1}$ (the engine divides the propensity by $M$ once), and zeroth-order sources in $\mu\mathrm{M\,s^{-1}}$ per voxel.
Rate constants are effective values of the model, not always the microscopic constants of a single enzyme; where a value is a lumped or rescaled quantity this is stated.

\paragraph{Two numerical conventions that affect propensities.}
The engine evaluates the raw product of the educt counts, so that $2\mathrm X$ would react with itself at rate $\propto n^2$ instead of $n(n-1)$.
At the copy numbers of the excitable element (a few molecules per cell at rest) this is a 23\,\% error, and it acts as a spurious first-order source when molecules are spread over 13 voxels.
Same-species reactions are therefore written with one educt and a falling-factorial factor (Table~\ref{tab:m0rxn}), which makes the per-voxel decomposition exact and the 13-voxel cell equivalent to one well-mixed compartment.
Second, Hill gates with $n\ge6$ are written in the reciprocal form $[1+(K/x)^n]^{-1}$, because the propensity arithmetic is single precision and $x^8$ overflows for $x>65\,536$, which would make the gate undefined.

\paragraph{Provenance of the constants.}
Every constant carries a basis tag in the constants tables:
\textbf{L}~a literature value used directly;
\textbf{D}~derived from a literature anchor by a stated relation;
\textbf{C}~calibrated on the model against a stated target behaviour, which is not an independent measurement;
\textbf{E}~an estimate or design choice for which no measurement is available.
Tags~C and E are the ones a reader should treat with caution.
The number of C/E constants is large by necessity: many molecular abundances and rates of this pathway have never been measured, and the calibrations aim at reproducing the measured \emph{behaviour} (kinetics, thresholds, periods) of the cell rather than individual molecular values.

\paragraph{Constants that are deliberately not literature values.}
Three departures are made explicit because they trade realism for signal-to-noise or tractability:
(i)~the G-protein to receptor ratio is 60 instead of about 3 (Module~2) to raise molecule counts in the gradient-sensing stages;
(ii)~the uniform acceleration of all Module-4 rates by a factor of 20, which compresses time scales that were not measured;
(iii)~the PdsA--PdiA dissociation constant is $10^4$ times weaker than measured (Module~11).

\subsection{Hill gates}\label{sec:hill}

Hill functions $x^n/(K^n+x^n)$ appear in the network wherever a lumped, thresholded step replaces a multi-step process that is not resolved by the model.
They are not used to claim molecular cooperativity of $n$ binding events.
A gate is justified when (i)~a switch-like transition is required for the function of the module, (ii)~its half-point $K$ can be placed relative to the resting and stimulated values of its input, and (iii)~its steepness $n$ is the smallest value that achieves the required contrast.
Table~\ref{tab:hill} lists all gates of the default network, with the value of the gate at rest and at the stimulated level of its input where these are known from the model.
Two further rules were enforced when placing the half-points: the activated fraction of a switch must be neither near zero nor near one at its working point (a saturated gate cannot report a difference), and each $K$ is placed between the resting and the driven value of its input.

\begin{table*}[!t]
\centering
\caption{Hill gates and saturable kinetics of the default network. $x$ is the input of the gate.}
\label{tab:hill}
\footnotesize
\renewcommand{\arraystretch}{1.2}
\begin{tabularx}{\textwidth}{@{}l l l c l >{\raggedright\arraybackslash}X@{}}
\toprule
Gate & Reaction & Form & $n$ & $K$ & Justification \\
\midrule
Competence $c(H)$ & 0.9a & $H^n/(K^n+H^n)$ & 4 & $60$ molec/voxel & Developmental switch between non-firing and firing cells; $c(100)=0.885$. Phenomenological (Hunger is an abstract clock), steep to give a sharp onset. \\
Refractory block $b(\mathrm{Rf})$ & 0.9a, 8.14 & $K^n/(K^n+x^n)$ & 4 & $450$ molec/cell & Decreasing gate that must shut the seed and the relay hard after a pulse and reopen within minutes. $K$ and $\tau_R$ were set on the measured recovery curve (blocked at 300\,s, open by 420\,s); $n=4$ makes the boundary sharp. \\
ACA gate & \mkxx{5}{1}{3}{0.9f} & $x^n/(K^n+x^n)$ & 4 & $192$ molec/voxel ERK2$^*$ & Opens the cyclase only after ERK2$^*$ has disarmed RegA (ERK2 was first reported to be \mkt{1}{needed for ACA activation}~\cite{segall1995}, but \mkt{3}{is not essential for it}~\cite{maeda2004}, and \mkt{7}{controls the phosphodiesterase RegA}~\cite{maeda2004,laub1998}); enforces the order ERK2$\to$ACA. \\
PTEN eviction & 4.6c & $x^n/(K^n+x^n)$ & 2 & $1.31\times10^5$ PIP3/voxel & Realises mutual inhibition PIP3 $\dashv$ PTEN~\cite{matsuoka2018}; $n=2$ is the minimum cooperativity for winner-take-all behaviour. $K$ (17\,\% of the lipid pool) sits on the steep flank; beyond $\approx\!19.5\,\%$ the domain becomes an orientation-independent latch. \\
PIP3 relay gate $g$ & 8.14 & $[1+(K/x)^n]^{-1}$ & 8 & $5000$ PIP3/voxel & The relay must open at the PIP3 evoked by 2\,nM cAMP ($\approx\!5000$/voxel; gate $0.51$) and stay shut at the resting level ($\approx\!10^3$; gate $<\!10^{-5}$), a contrast of only about 5. The exponent $n=8$ was set when the peak-to-plateau ratio was 2.7, for which a 300-fold suppression of the standing drive needs $n\approx8$, and was kept when $K$ was re-anchored. It stands for the lumped chain PIP3$\to$PKB$\to$CRAC$\to$cyclase priming, not one binding event. Reciprocal form (Sec.~\ref{sec:conv}). \\
NF1 saturable GAP & 3.7z & $x/(K+x)$ & 1 & $0.2\,\mu$M & Michaelis--Menten hydrolysis with $K\ll[\mathrm{Ras}]_\mathrm{tot}$ (zero-order regime), which produces rectification of Ras activation~\cite{nakajima2014}; $K$ from~\cite{wiesmuller1992,coyle2016}. \\
PdsA induction $P(H)$ & 11.D1b, 11.E1 & $H^n/(K^n+H^n)$ & 2 & $25$ molec/voxel & Starvation gate of the aggregation-stage \emph{pdsA}/\emph{pdiA} expression; opens early relative to the Hunger ramp, $P(100)=0.94$. \\
PdiA repression $r$ & 11.E1 & $[1+(x/K)^n]^{-1}$ & 2 & $0.05\,\mu$M cAMP$_s$ & Repression of PdiA synthesis by the slowly filtered cAMP signal~\cite{yeh1978}; $n=2$ gives a graded but steeper-than-hyperbolic switch. \\
Memory writing (optional) & M.1 & $x^n/(K^n+x^n)$ & 3 & $5000$ PIP3/voxel & Writes the mark only on the PIP3 rise ($H(\text{rest})\approx0.03$, $H(7500)\approx0.77$). \\
\bottomrule
\end{tabularx}
\end{table*}

\subsection{Reproducibility}
All tables below are generated from the values that the simulation loads by default; every constant can be overridden from the environment for sensitivity studies.
Two configurations of the code (a two-state and a multi-affinity receptor, Module~1) exist; only the two-state receptor is used.

\subsection{Module 0: Developmental State, Excitable Element and Refractory Clock}\label{sec:m0}

\paragraph{Role.}
Before collective waves start, isolated starved cells emit cAMP pulses spontaneously, on average \mk{1}{about once per 12\,min}~\cite{ford2023} and \mk{2}{at most about once per 4\,min}~\cite{ford2023}, with inter-pulse intervals that \mk{3}{follow a gamma distribution (shape 6.6, rate {$1.25\,\mathrm{min^{-1}}$})}~\cite{ford2023}, i.e.\ \mk{5}{a mean of 5.3\,min}~\cite{ford2023}.
Such \mk{9}{noise-driven pulsing of single cells below the threshold for collective behaviour is a key step in the onset of oscillations}~\cite{gregor2010,sgro2015}.
The molecular circuit that generates periodic cAMP in populations couples \mk{5}{cAMP-induced activation of ACA and ERK2}~\cite{laub1998,maeda2004}, \mk{6}{inhibition of ERK2 by PKA}~\cite{laub1998,maeda2004} and \mk{7}{phosphorylation of the phosphodiesterase RegA by ERK2}~\cite{laub1998,maeda2004}, and produces a period of \mk{8}{about 7\,min}~\cite{laub1998,maeda2004}.
Module~0 replaces the unresolved part of this circuit by a compact excitable element and adds the two variables that the collective behaviour needs: a developmental state and a refractory clock.

\paragraph{Why the driver is X and not cAMP.}
In the literature circuit cAMP is the input that activates ACA and ERK2.
In the reactions of this module (Table~\ref{tab:m0rxn}) ERK2 and ACA are instead activated by the abstract activator X (0.9e, 0.9f); this is a deliberate design choice and not a measured dependence.
Wiring cAMP into these steps does not work in this network, for two reasons found while building the model.
First, every route from the receptor to the cyclase passes through receptor occupancy, so an isolated cell without external cAMP has no input to ACA at all and can never fire, although isolated starved cells do pulse spontaneously~\cite{ford2023}.
Second, a feedback of intracellular cAMP onto ACA cannot produce a pulse: all paths from cAMP$_i$ to ACA run through PKA, whose edges are inhibitory, and an added cooperative autocatalysis on cAMP$_i$ either settles on a plateau or latches in the high state; neither has a reset.
A pulse needs a fast autocatalysis together with a slow recovery, and the receptor-independent arm of the network has neither.
X supplies both (Schl\"ogl autocatalysis with the recovery variable Y, a FitzHugh--Nagumo pair), is seeded by the Hunger-gated ignition 0.9a so that it needs no receptor input, and is therefore the node at which an isolated cell decides to fire.
X thus stands for the unresolved part of the circuit between cAMP and ACA/ERK2 and is not a molecule that has been identified.
Consequently the model contains no direct activation of ACA or ERK2 by cAMP in the table of this module; the cAMP-dependence of the circuit enters only through the relay that fires X (Module~8).

\begin{table*}[!t]
\centering
\caption{Module~0 reactions: Hunger, the excitable element (X, Y), the refractory clock (Rf) and the ERK2$\to$ACA output stage. For the species of the excitable element $N_s$ is the copy number of species $s$ \emph{per cell}; $c(\mathrm H)$ and $b(\mathrm{Rf})$ are the Hill gates of Sec.~\ref{sec:hill}.}
\label{tab:m0rxn}
\footnotesize
\renewcommand{\arraystretch}{1.15}
\begin{tabularx}{\textwidth}{@{}l>{\raggedright\arraybackslash}p{0.37\textwidth}>{\raggedright\arraybackslash}X@{}}
\toprule
ID & Reaction & Propensity (rate law)\\
\midrule
0.8a & $\varnothing \to \mathrm{H}$ & $k_H\,H_{\max}\ \text{(zeroth order)}$ \\
0.8b & $\mathrm{H} \to \varnothing$ & $k_H\, N_\mathrm{H}$ \\
0.9a & $\varnothing \to \mathrm{X}$ & $k_0\, c(\mathrm{H})\, b(\mathrm{Rf})\quad[\text{molec\,s}^{-1}\text{cell}^{-1}]$ \\
0.9b & $2\mathrm{X} \to 3\mathrm{X}$ & $\tfrac{k_1}{2}N_\mathrm{X}(N_\mathrm{X}-1)$ \\
0.9c & $3\mathrm{X} \to 2\mathrm{X}$ & $\tfrac{k_6}{6}N_\mathrm{X}(N_\mathrm{X}-1)(N_\mathrm{X}-2)$ \\
0.9d & $\mathrm{X} \to \varnothing$ & $k_2 N_\mathrm{X}$ \\
0.9y1 & $2\mathrm{X} \to 2\mathrm{X}+\mathrm{Y}$ & $\tfrac{k_3}{2}N_\mathrm{X}(N_\mathrm{X}-1)$ \\
0.9y2 & $\mathrm{X}+\mathrm{Y} \to \mathrm{Y}$ & $k_5 N_\mathrm{X}N_\mathrm{Y}$ \\
0.9y3 & $\mathrm{Y} \to \varnothing$ & $k_4 N_\mathrm{Y}$ \\
0.9r1 & $2\mathrm{X} \to 2\mathrm{X}+\mathrm{Rf}$ & $\tfrac{k_{Rp}}{2}N_\mathrm{X}(N_\mathrm{X}-1)$ \\
0.9r2 & $\mathrm{Rf} \to \varnothing$ & $k_{Rd} N_\mathrm{Rf}$ \\
0.9e & $2\mathrm{X}+\mathrm{ERK2} \to 2\mathrm{X}+\mathrm{ERK2}^*$ & $k_{XE}\,[\mathrm{X}]^2[\mathrm{ERK2}]\ \text{(falling factorial in }N_\mathrm{X})$ \\
0.9e$'$ & $\mathrm{ERK2}^* \to \mathrm{ERK2}$ & $k_{Ed}[\mathrm{ERK2}^*]$ \\
0.9f & $\mathrm{X}+\mathrm{ACA} \to \mathrm{X}+\mathrm{ACA}^*$ & $k_{XA}[\mathrm{X}][\mathrm{ACA}]\,\dfrac{[\mathrm{ERK2}^*]^4}{K_E^4+[\mathrm{ERK2}^*]^4}$ \\
\bottomrule
\end{tabularx}
\end{table*}

\begin{table*}[!t]
\centering
\caption{Module~0 constants. Basis: L literature value, D derived from a literature anchor, C calibrated on the model against a stated target, E estimate or design choice (Sec.~\ref{sec:conv}).}
\label{tab:m0const}
\footnotesize
\renewcommand{\arraystretch}{1.15}
\begin{tabularx}{\textwidth}{@{}l r l c >{\raggedright\arraybackslash}X@{}}
\toprule
Symbol & Value & Unit & Basis & Meaning, justification, source\\
\midrule
$H_{\max}$ & \num{100} & $\mathrm{molec/voxel}$ & E & Steady-state level of the abstract \emph{Hunger} variable; held at full developmental competence. \\
$k_H$ & \num{2.8e-4} & $\mathrm{s^{-1}}$ & E & Relaxation of Hunger ($\tau\approx 1$\,h); the production rate is $H_{\max}k_H$. A slow developmental clock, not a measured rate. \\
$H_{1/2},\ n_c$ & \num{60},\ 4 & $\mathrm{molec/voxel}$ & E & Half-competence level and steepness of the competence gate; $c(H_{\max})=0.885$. \\
$k_0$ & \num{3.5} & $\mathrm{molec\,s^{-1}\,cell^{-1}}$ & C & Ignition (seed) rate at full competence. Set on the stochastic engine so that the mean inter-pulse interval of an isolated cell is 726\,s; target \mkt{1}{{$\approx\!12$}\,min}~\cite{ford2023}. \\
$k_1$ & \num{0.08} & $\mathrm{molec^{-1}s^{-1}}$ & C & Autocatalysis ($2\mathrm X\to3\mathrm X$); creates the excitation threshold ($\approx\!23$ molecules). \\
$k_6$ & \num{5e-4} & $\mathrm{molec^{-2}s^{-1}}$ & C & Cubic self-limitation; bounds the spike at $\approx\!456$ molecules per cell. \\
$k_2$ & \num{1} & $\mathrm{s^{-1}}$ & C & Basal removal of X (relaxation time 1\,s). \\
$k_3$ & \num{6.25e-5} & $\mathrm{molec^{-1}s^{-1}}$ & C & Spike-triggered production of the recovery variable Y. \\
$k_5$ & \num{0.05} & $\mathrm{molec^{-1}s^{-1}}$ & C & Y-mediated spike termination. \\
$k_4$ & \num{0.01} & $\mathrm{s^{-1}}$ & C & Decay of Y ($\tau=100$\,s); sets the recovery of the element. \\
$k_{Rp},\ k_{Rd}$ & \num{5e-3},\ \num{5e-3} & $\mathrm{molec^{-1}s^{-1}},\ \mathrm{s^{-1}}$ & C & Charging and decay ($\tau_R=200$\,s) of the refractory clock. \\
$K_R,\ n_R$ & \num{450},\ 4 & $\mathrm{molec/cell}$ & C & Half-block level and steepness of the refractory gate (Sec.~\ref{sec:hill}). \\
$k_{XE}$ & \num{1823} & $\mu\mathrm{M}^{-2}\mathrm{s}^{-1}$ & C & X$^2$-driven activation of ERK2 (cell units); second order in X to raise the contrast between spike and resting leak. \\
$k_{Ed}$ & \num{0.02} & $\mathrm{s^{-1}}$ & E & ERK2 dephosphorylation ($\tau=50$\,s); ERK2 activation is \mkt{1}{transient}~\cite{maeda1996}. \\
$k_{XA}$ & \num{65} & $\mu\mathrm{M}^{-1}\mathrm{s}^{-1}$ & C & X-driven ACA activation, gated on ERK2$^*$; sets the $\approx\!10$\,s rise of ACA$^*$. \\
$K_E,\ n_E$ & \num{192},\ 4 & $\mathrm{molec/voxel}$ & C & ERK2$^*$ level (12.5\,\% of the pool) at half-open ACA gate; enforces the ordering ERK2$\to$ACA. \\
\midrule
$N_\mathrm{vox}$ & \num{13} &  & E & Voxels per cell (plus-shaped footprint of $52\,\mu\mathrm{m}^2$, i.e.\ about 8\,$\mu$m diameter, in a 2\,$\mu$m thick layer, $104\,\mu\mathrm{m}^3$); real cells have a \mkt{4}{diameter of about 10\,{$\mu$}m}~\cite{ford2023}, so the cell volume is underestimated. $M_c=13\times4817.6=62\,629$ molecules\,$\mu\mathrm{M}^{-1}$ per cell. \\
$[\mathrm{ERK2}]_\mathrm{tot}$ & \num{1540} & $\mathrm{molec/voxel}$ & E & Pool size on the scale of the other kinases ($0.32\,\mu$M). \\
$[\mathrm{ACA}]_\mathrm{tot}$ & \num{865} & $\mathrm{molec/voxel}$ & C & Chosen so that one pulse exports $\approx\!10^7$ cAMP molecules per cell at $k_9+k_{11}=20\,\mathrm{s^{-1}}$ (Module~9). \\
$D_\mathrm{X},D_\mathrm{Y},D_\mathrm{Rf},D_\mathrm{ERK2}$ & \num{10} & $\mu\mathrm{m^2/s}$ & E & Generic cytosolic value; the cell is mixed in $\approx\!2.5$\,s, which matters because the gates are evaluated per voxel. \\
\bottomrule
\end{tabularx}
\end{table*}


\paragraph{Reactions.}
\begin{description}\itemsep0pt
\item[0.8a/b] An abstract \emph{Hunger} variable $H$ represents the developmental state of a starved cell; it is not a molecule. It scales the competence to fire (gate $c(H)$), the cAMP synthesis rate (9.1) and the induction of the extracellular phosphodiesterase system (Module~11). In the simulations $H$ starts at its steady state $H_{\max}$ and fluctuates around it, so it acts as an approximately constant developmental level.
\item[0.9a] Spontaneous ignition: a rare, competence-gated seed that injects X, the activator, at $k_0=3.5$ molecules\,s$^{-1}$ per cell. It is switched off by the refractory clock, which makes the block hard: with the seed shut, X decays through 0.9d and the autocatalysis 0.9b, which needs $N_\mathrm{X}\ge2$, has nothing to amplify.
\item[0.9b--d] The autocatalytic step 0.9b creates an excitation threshold, the cubic loss 0.9c bounds the spike and keeps the stochastic simulation finite, and the linear loss 0.9d sets the resting state. These three reactions are the \mk{4}{chemical Schl\"ogl model}~\cite{schlogl1972} of a \mk{3}{bistable}~\cite{schlogl1972}/excitable system, in a form in which the buffered reservoirs A, B and C of the original scheme (A${}+2$X$\rightleftharpoons3$X, B${}+$X$\rightleftharpoons$C) are absorbed into the rate constants: 0.9b is the forward and 0.9c the reverse direction of A${}+2$X$\to3$X (the $k_1an^2-k_1'n^3$ of the original rate law), and 0.9d is B${}+$X$\to$C. The reverse direction C$\to$B${}+$X is not a separate reaction; its role as the source of X is taken by the gated ignition 0.9a. The propensities are the stochastic (falling-factorial) form of the Schl\"ogl rate law.
\item[0.9y1--3] The recovery variable Y is produced only by spikes (quadratic in X, so resting fluctuations do not raise it), terminates the spike by removing X (0.9y2) and decays slowly (0.9y3). X and Y form the chemical counterpart of the \mk{2}{FitzHugh--Nagumo excitable system with activator and delayed inhibitor}~\cite{fitzhugh1961,nagumo1962}. With the constants of Table~\ref{tab:m0const} and the competence gate at $H_{\max}$ ($c=0.885$), the cell-level nullcline gives a resting level $\approx3.4$, an excitation threshold $\approx24$ and a spike maximum $\approx456$ molecules per cell.
\item[0.9r1/2] A refractory clock Rf is charged by spikes of X (quadratic, as for Y) and decays with $\tau_R=200$\,s. Y cannot serve as the clock because it also terminates the spike, and its time constant would then be tied to the pulse shape. Charging the clock from X makes the block \emph{source independent}: a cell that has fired, spontaneously or after a relayed signal (Module~8), cannot be fired again by either route until the clock has decayed. The model requires a refractory period of several minutes, commensurate with the \mk{8}{$\approx$7\,min oscillation period}~\cite{maeda2004} and the \mk{2}{$\approx$4\,min minimum interval between pulses of single cells}~\cite{ford2023}.
\item[0.9e] X activates ERK2 in a second-order step. In the cell, \mk{5}{ERK2 is activated by cAMP}~\cite{maeda1996,laub1998,maeda2004} through the receptor cAR1; X takes the place of this cAMP signal for the reason given above (an isolated cell has no external cAMP, so a receptor-driven step could never start a pulse). ERK2$^*$ is deliberately long-lived, so it integrates X; in the calibration a first-order drive left ERK2$^*$ at 5\,\% of the pool at rest, which kept the RegA brake 87\,\% off before any spike. The second-order dependence squares the spike-to-leak contrast and can be read as a lumped ultrasensitive step; the second order in X has no direct experimental counterpart.
\item[0.9e$'$] First-order return of ERK2$^*$ to ERK2. No dephosphorylation rate or phosphatase has been measured; the rate is chosen to reproduce the measured time course of ERK2 after a cAMP stimulus, \mk{1}{a peak at 50--60\,s and a return to near-basal activity within minutes}~\cite{maeda1996,laub1998,aubry1997} (4\,min in~\cite{maeda1996}, 5--8\,min in~\cite{aubry1997}). In the Laub--Loomis circuit the only measured-type ERK2 inactivation is by PKA~\cite{laub1998} (reaction 10.D3); a PKA-independent decay is not part of it, so 0.9e$'$ is an added term. The value $\tau=50$\,s ($>\!95\,\%$ decayed after 3\,min) lies at the fast end of the reported range and was fixed by the requirement that resting ERK2$^*$ stays low.
\item[0.9f] X activates the adenylyl cyclase ACA once ERK2$^*$ has risen. ERK2 was reported to be \mk{1}{required for receptor-mediated activation of adenylyl cyclase}~\cite{segall1995}, but later work showed that \mk{3}{ERK2 is not essential for ACA activation and acts on cAMP accumulation through RegA}~\cite{maeda2004}; the ERK2$^*$ gate on ACA is therefore a design choice that orders the two events, not a measured dependence. ERK2$^*$ disarms RegA (Module~10). The ERK2$^*$ gate makes ACA follow the release of the RegA brake by a few seconds, so the cAMP made by ACA$^*$ is not destroyed immediately.
\end{description}

\paragraph{Constants.}
The parameters of 0.9a--d and 0.9y are calibrated, not measured: the excitable core has three targets (excitation threshold, spike height and recovery time), and the values in Table~\ref{tab:m0const} reproduce them.
The ignition rate $k_0$ is the only parameter that sets the firing rate.
It was calibrated on the stochastic engine, because neither a well-mixed reduction (about $2\times$ too slow: ignition nucleates in one voxel before diffusion levels it out) nor a 13-voxel deterministic reduction predicts it.
With the refractory clock active, $k_0=3.5$ molecules\,s$^{-1}$ gives a mean interval of 726\,s (95\,\% interval 706--747\,s, 8 seeds), close to the \mk{1}{12\,min} of~\cite{ford2023}; the previous value 2.52 gave 930\,s because the block adds about 390\,s of dead time.
The 12\,min of the text of~\cite{ford2023} and the mean of 5.3\,min of its fitted gamma distribution (Fig.~1E) differ; the model follows the former. \emph{Limitation:} because the block is deterministic and makes up a large part of the interval, the coefficient of variation of the isolated pacemaker is 0.16, below the value $1/\sqrt{6.6}\approx0.39$ implied by the \mk{3}{gamma distribution} measured in~\cite{ford2023}.
In populations the period converges onto the block; for 1000 cells it is $467\pm16$\,s (7.8\,min), of the order of the \mk{8}{7\,min period} of~\cite{maeda2004}.
The clock parameters $k_{Rp},k_{Rd},K_R$ were set on the measured recovery curve of the relay (no relayed pulse within 300\,s of the previous pulse, one in three at 360\,s, all by 420\,s; three seeds), not from the analytical estimate, because the block lifts when the gated relay drive crosses the excitation threshold, which depends on the relay gain.

The Hill gates $c(H)$, $b(\mathrm{Rf})$ and the ERK2$^*$ gate are justified in Table~\ref{tab:hill}.
The choice $n=4$ for the three is not a cooperativity claim; it makes the transition sharp enough that the gate is either nearly open or nearly shut at the values the input actually takes.

\subsection{Module 1: cAMP Receptor Binding and Desensitisation}\label{sec:m1}

\paragraph{Role.}
cAMP is sensed by the G-protein-coupled receptor cAR1.
Binding to intact aggregation-competent cells is heterogeneous: \mk{M1-slow4}{about 4\,\% of the sites dissociate slowly ($K_d=12.5$\,nM)}~\cite{vanhaastert1984}, \mk{M1-fast-high}{about 40\,\% are fast, high-affinity sites ($K_d=60$\,nM)}~\cite{vanhaastert1984} and \mk{M1-fast-low}{about 60\,\% are fast, low-affinity sites ($K_d=450$\,nM)}~\cite{vanhaastert1984}; \mk{M1-conv9}{under cAMP the high-affinity sites convert to low-affinity ones with first-order kinetics ($t_{1/2}\approx9$\,s)}~\cite{vanhaastert1984}.
Occupancy \mk{M1-ser-phos}{induces phosphorylation of serine clusters in the C-terminal domain}~\cite{vaughan1988,hereld1994}, which is \mk{M1-llb-corr}{correlated with the loss of ligand binding}~\cite{caterina1995a,caterina1995b} and \mk{M1-aff-red}{reduces the affinity of the receptor rather than removing it from the surface}~\cite{caterina1995a,caterina1995b}.
The model resolves the two fast classes, and represents the low-affinity class as the phosphorylated receptor; the slow class (4\,\% of sites) is omitted.
Each receptor form has the binding topology of the \mk{M1-bis-red}{reduced-order receptor model} of Biswas et al.~\cite{biswas2021} (Table~1, rows 27/28), $\mathrm R+\mathrm{cAMP}\rightleftharpoons\mathrm{RC}$, to which the model adds a phosphorylated pair; the \mk{M1-bis-full}{full three-class model} of~\cite{biswas2021} is available as an option (Table~\ref{tab:m1brx}) and is not used here.

\begin{table*}[!t]
\centering
\caption{Module~1 reactions (default two-state receptor). R is the unphosphorylated cAR1 receptor, R$_p$ its phosphorylated form, RC and R$_p$C the cAMP-bound states; A is the cytosolic adapter of reactions 1.8. A coloured reaction is the one to which the equally coloured statement in the text and the equally coloured passage in the cited paper refer; a small square appends a further statement.}
\label{tab:m1rxn}
\footnotesize
\renewcommand{\arraystretch}{1.15}
\begin{tabularx}{\textwidth}{@{}l>{\raggedright\arraybackslash}p{0.37\textwidth}>{\raggedright\arraybackslash}X@{}}
\toprule
ID & Reaction & Propensity (rate law)\\
\midrule
\mkn{M1-bis-red}{M1-kon-tab,M1-fast-high}{1.1} & \mkn{M1-bis-red}{M1-kon-tab,M1-fast-high}{$\mathrm{R}+\mathrm{cAMP}_e \to \mathrm{RC}$} & $k_\mathrm{on}[\mathrm{R}][\mathrm{cAMP}_e]$ \\
\mkn{M1-bis-red}{M1-kon-tab,M1-fast-high,M1-koff1s}{1.2} & \mkn{M1-bis-red}{M1-kon-tab,M1-fast-high,M1-koff1s}{$\mathrm{RC} \to \mathrm{R}+\mathrm{cAMP}_e$} & $k_\mathrm{off}[\mathrm{RC}]$ \\
\mkn{M1-fast-low}{M1-aff-red,M1-aff35}{1.3} & \mkn{M1-fast-low}{M1-aff-red,M1-aff35}{$\mathrm{R}_p+\mathrm{cAMP}_e \to \mathrm{R}_p\mathrm{C}$} & $k_{\mathrm{on},p}[\mathrm{R}_p][\mathrm{cAMP}_e]$ \\
\mkn{M1-fast-low}{M1-aff-red,M1-aff35,M1-koff1s}{1.4} & \mkn{M1-fast-low}{M1-aff-red,M1-aff35,M1-koff1s}{$\mathrm{R}_p\mathrm{C} \to \mathrm{R}_p+\mathrm{cAMP}_e$} & $k_{\mathrm{off},p}[\mathrm{R}_p\mathrm{C}]$ \\
\mkn{M1-phos45}{M1-ser-phos,M1-llb-corr,M1-conv9}{1.5} & \mkn{M1-phos45}{M1-ser-phos,M1-llb-corr,M1-conv9}{$\mathrm{RC} \to \mathrm{R}_p\mathrm{C}$} & $k_\mathrm{ph}[\mathrm{RC}]$ \\
\mkt{M1-basal}{1.5b} & \mkt{M1-basal}{$\mathrm{R} \to \mathrm{R}_p$} & $k_{\mathrm{ph}0}[\mathrm{R}]$ \\
\mkn{M1-loss80}{M1-h10}{1.6} & \mkn{M1-loss80}{M1-h10}{$\mathrm{R}_p\mathrm{C} \to \mathrm{RC}$} & $k_\mathrm{dp}[\mathrm{R}_p\mathrm{C}]$ \\
\mkt{M1-deph2}{1.7} & \mkt{M1-deph2}{$\mathrm{R}_p \to \mathrm{R}$} & $k_\mathrm{b}[\mathrm{R}_p]$ \\
\mkt{M1-fcd}{1.8a} & \mkt{M1-fcd}{$\mathrm{RC}+\mathrm{A} \to \mathrm{RC}+\mathrm{A}^*$} & $k_{A}[\mathrm{RC}][\mathrm{A}]$ \\
\mkt{M1-fcd}{1.8b} & \mkt{M1-fcd}{$\mathrm{A}^* \to \mathrm{A}$} & $k_{A,\mathrm{off}}[\mathrm{A}^*]$ \\
\bottomrule
\end{tabularx}
\end{table*}

\begin{table*}[!t]
\centering
\caption{Module~1 constants (default receptor model).}
\label{tab:m1const}
\footnotesize
\renewcommand{\arraystretch}{1.15}
\begin{tabularx}{\textwidth}{@{}l r l c >{\raggedright\arraybackslash}X@{}}
\toprule
Symbol & Value & Unit & Basis & Meaning, justification, source\\
\midrule
$R_\mathrm{tot}$ & \num{3077} & $\mathrm{molec/voxel}$ & E & 40\,001 receptors per cell ($0.64\,\mu$M). Cells developed for 4\,h carry about $7\times10^4$ binding sites~\cite{johnson1991}, and Biswas et al.\ use $7\times10^4\pm5\times10^3$~\cite{biswas2021}; the model value is a lower design choice. \\
$k_\mathrm{on}$ & \num{7.5} & $\mu\mathrm{M}^{-1}\mathrm{s}^{-1}$ & L & Fast high-affinity class H of Ref.~\cite{biswas2021}, Table~1, taken from~\cite{vanhaastert1984}. \\
$k_\mathrm{off}$ & \num{0.45} & $\mathrm{s^{-1}}$ & L & Gives $K_d=k_\mathrm{off}/k_\mathrm{on}=60$\,nM, the fast high-affinity class ($\approx\!40$\,\% of sites) of~\cite{vanhaastert1984}. \\
$\varepsilon_p$ & \num{0.125} &  & C & Ratio of on-rates of R$_p$ and R. Chosen so that $K_{d,p}=480$\,nM, close to the 450\,nM low-affinity class of~\cite{vanhaastert1984}; the loss of affinity on phosphorylation is documented in~\cite{caterina1995b,caterina1995a}. \\
$k_{\mathrm{on},p}$ & \num{0.938} & $\mu\mathrm{M}^{-1}\mathrm{s}^{-1}$ & D & $\varepsilon_p k_\mathrm{on}$. The affinity loss is carried by the on-rate, so R$_p$ still binds and releases cAMP. \\
$k_{\mathrm{off},p}$ & \num{0.45} & $\mathrm{s^{-1}}$ & L & Equal to $k_\mathrm{off}$; the fast dissociation half-life of $\approx\!1$\,s is common to both fast classes~\cite{vanhaastert1984}. \\
$k_\mathrm{ph}$ & \num{0.02} & $\mathrm{s^{-1}}$ & L & Agonist-induced phosphorylation ($t_{1/2}=35$\,s) against $t_{1/2}\approx45$\,s at saturating cAMP in~\cite{vaughan1988}. \\
$k_{\mathrm{ph}0}$ & \num{5e-4} & $\mathrm{s^{-1}}$ & E & Ligand-independent phosphorylation (40$\times$ slower); a basal component exists~\cite{vaughan1988,hereld1994}. \\
$k_\mathrm{dp}$ & \num{5e-3} & $\mathrm{s^{-1}}$ & C & Dephosphorylation of the bound receptor ($t_{1/2}=139$\,s). Sets the phosphorylated fraction $k_\mathrm{ph}/(k_\mathrm{ph}+k_\mathrm{dp})=80\,\%$ under sustained saturating cAMP, matching the $>\!80\,\%$ loss of high-affinity binding~\cite{vanhaastert1984,johnson1991}. \\
$k_\mathrm{b}$ & \num{5.78e-3} & $\mathrm{s^{-1}}$ & L & $t_{1/2}=2$\,min, the receptor dephosphorylation half-time after removal of the stimulus~\cite{vaughan1988}. \\
$k_A,\ k_{A,\mathrm{off}}$ & \num{0.07},\ \num{0.02} & $\mu\mathrm{M}^{-1}\mathrm{s}^{-1},\ \mathrm{s^{-1}}$ & E & Adapter charged by RC ($\tau=50$\,s); dormant in the default network (see text). \\
$[\mathrm{A}]_\mathrm{tot}$ & \num{154} & $\mathrm{molec/voxel}$ & E & Adapter pool. \\
$D_{\mathrm{R}}$ & \num{0.024} & $\mu\mathrm{m^2/s}$ & L & Common lateral diffusion constant of transmembrane proteins in \emph{Dictyostelium}, $0.024\pm0.004\,\mu\mathrm{m}^2/\mathrm{s}$ ($n=27$)~\cite{takebayashi2023}; Biswas et al.\ use $0.027$~\cite{biswas2021,ueda2001}. \\
$D_{\mathrm{A}}$ & \num{10} & $\mu\mathrm{m^2/s}$ & E & Generic cytosolic value. \\
\bottomrule
\end{tabularx}
\end{table*}

\begin{table*}[!t]
\centering
\caption{Optional full multi-affinity receptor module of Biswas et al.~\cite{biswas2021} (Table~1, rows 1--26), selectable in the code but \emph{not} used in the results of this paper. Classes H, L, S: high, low and slow affinity; P is the free-phosphate pool; $^P$ marks phosphorylated states. Each row lists the forward and reverse constants.}
\label{tab:m1brx}
\footnotesize
\renewcommand{\arraystretch}{1.15}
\begin{tabularx}{\textwidth}{@{}l>{\raggedright\arraybackslash}p{0.37\textwidth}>{\raggedright\arraybackslash}X@{}}
\toprule
ID & Reaction & Propensity (rate law)\\
\midrule
1/2 & $\mathrm{H}+\mathrm{cAMP}\rightleftharpoons \mathrm{H{:}C}$ & $k_H=\num{7.5},\ k_{-H}=\num{0.45}$ \\
3/4 & $\mathrm{L}+\mathrm{cAMP}\rightleftharpoons \mathrm{L{:}C}$ & $k_L=\num{2.2},\ k_{-L}=\num{1}$ \\
5/6 & $\mathrm{S}+\mathrm{cAMP}\rightleftharpoons \mathrm{S{:}C}$ & $k_S=\num{4},\ k_{-S}=\num{0.05}$ \\
7/8 & $\mathrm{H}\rightleftharpoons \mathrm{L}$ & $k_{HL}=\num{0.08},\ k_{-HL}=\num{0.0536}$ \\
9/10 & $\mathrm{H{:}C}\rightleftharpoons \mathrm{L{:}C}$ & $k_{HLC}=\num{0.08},\ k_{-HLC}=\num{7.1e-3}$ \\
11/12 & $\mathrm{H{:}C}+\mathrm{P}\rightleftharpoons {}^P\mathrm{H{:}C}$ & $k_{HCP}=\num{5.3e-4},\ k_{-HCP}=\num{8e-4}$ \\
13/14 & ${}^P\mathrm{H}+\mathrm{cAMP}\rightleftharpoons {}^P\mathrm{H{:}C}$ & $k_{PHC}=\num{0.04},\ k_{-PHC}=\num{2.42e-3}$ \\
15/16 & $\mathrm{H}+\mathrm{P}\rightleftharpoons {}^P\mathrm{H}$ & $k_{PH}=\num{5.3e-4},\ k_{-PH}=\num{8e-4}$ \\
17/18 & $\mathrm{L{:}C}+\mathrm{P}\rightleftharpoons {}^P\mathrm{L{:}C}$ & $k_{LCP}=\num{5.3e-4},\ k_{-LCP}=\num{8e-4}$ \\
19/20 & ${}^P\mathrm{L}+\mathrm{cAMP}\rightleftharpoons {}^P\mathrm{L{:}C}$ & $k_{PLC}=\num{5},\ k_{-PLC}=\num{1.5}$ \\
21/22 & $\mathrm{L}+\mathrm{P}\rightleftharpoons {}^P\mathrm{L}$ & $k_{PL}=\num{5.3e-4},\ k_{-PL}=\num{8e-4}$ \\
23/24 & ${}^P\mathrm{H}\rightleftharpoons {}^P\mathrm{L}$ & $k_{PHL}=\num{4.3e-4},\ k_{-PHL}=\num{2.9e-4}$ \\
25/26 & ${}^P\mathrm{H{:}C}\rightleftharpoons {}^P\mathrm{L{:}C}$ & $k_{PHLC}=\num{4.3e-4},\ k_{-PHLC}=\num{3.8e-5}$ \\
\bottomrule
\end{tabularx}
\end{table*}


\paragraph{Reactions.}
\begin{description}\itemsep0pt
\item[1.1/1.2] Reversible cAMP binding to the high-affinity receptor R, with $K_d=60$\,nM. This is the fastest step of the cascade ($\tau=1.1$\,s at $c=K_d$) and is close to quasi-equilibrium with everything downstream.
\item[1.3/1.4] Binding to the phosphorylated receptor R$_p$. Affinity is lost through the on-rate ($K_{d,p}=480$\,nM), so R$_p$ still binds and releases cAMP, i.e.\ it is slower to engage, not locked out. A low-affinity receptor is still on the steep part of its binding curve at concentrations where the high-affinity receptor is saturated; R$_p$ therefore extends the dynamic range of the stage.
\item[1.5/1.5b] Phosphorylation, agonist-induced (1.5) and basal (1.5b, 40$\times$ slower). Agonist-induced phosphorylation is measured with \mk{M1-phos45}{$t_{1/2}\approx45$\,s at saturating cAMP}~\cite{vaughan1988} and \mk{M1-basal}{a ligand-independent basal component exists}~\cite{vaughan1988}; the model uses $t_{1/2}=35$\,s.
\item[1.6/1.7] Dephosphorylation of the bound (1.6) and free (1.7) receptor. The free receptor is dephosphorylated with \mk{M1-deph2}{$t_{1/2}=2$\,min after removal of the stimulus}~\cite{vaughan1988}, so 1.7 sets the time to resensitise after a pulse. Under sustained saturating cAMP, 80\,\% of receptors are phosphorylated at steady state ($k_\mathrm{ph}/(k_\mathrm{ph}+k_\mathrm{dp})$), matching \mk{M1-loss80}{the loss of more than 80\,\% of the surface binding sites during prolonged stimulation}~\cite{johnson1991} and \mk{M1-h10}{the fall of the high-affinity fraction from about 40\,\% to about 10\,\%}~\cite{vanhaastert1984}.
\item[1.8a/b] A cytosolic adapter A is charged by occupied receptors and relaxes with $\tau=50$\,s. It was introduced as a fold-change-detecting gate for a direct receptor-to-cyclase edge (reaction (8.11)), motivated by \mk{M1-fcd}{the fold-change response of the cAMP relay}~\cite{kamino2017}. That edge is switched off ($k=0$), so the adapter has \emph{no downstream effect} in the default network; it is retained as an observable.
\end{description}

\paragraph{Constants.}
$k_\mathrm{on}$ and $k_\mathrm{off}$ are \mk{M1-kon-tab}{the high-affinity constants of Table~1 of}~\cite{biswas2021}, \mk{M1-kon-tab}{which in turn come from the analysis of}~\cite{vanhaastert1984}, and give $K_d=60$\,nM.
The on-rate ratio $\varepsilon_p=1/8$ makes $K_{d,p}=480$\,nM, close to the 450\,nM low-affinity class; the phosphorylation-induced affinity change was reported as \mk{M1-aff35}{3--5-fold for one low-affinity class} in the literature used by~\cite{biswas2021}, so the 8-fold ratio is at the upper end and \mk{M1-ratio75}{equals the ratio of the two fast classes}~\cite{vanhaastert1984}.
The receptor number, 40\,001 per cell, is lower than the \mk{M1-rtot}{$\approx7\times10^4$ sites measured in cells developed for 4\,h}~\cite{johnson1991} and used in~\cite{biswas2021}; it is a design choice and not a measurement.
The lateral diffusion constant $0.024\,\mu\mathrm{m^2/s}$ is \mk{M1-dif-t}{the mean over 27 transmembrane proteins measured in \emph{Dictyostelium}}~\cite{takebayashi2023}, close to \mk{M1-dif-u}{the value $0.027$ used by Biswas et al.}~\cite{biswas2021,ueda2001}.
A much faster phospho-cycle ($k_\mathrm{ph}=1$\,s$^{-1}$, about 50 times faster than the measured half-time of 45\,s) was tested as a way to de-saturate the receptor stage and rejected, because it made the unphosphorylated receptor resupply-limited and hence no longer a sensor; the present values are the literature-anchored ones.

\paragraph{Limitations.}
Ligand-induced phosphorylation is not necessary for adaptation of the downstream responses or for chemotaxis; \mk{M1-kim-ok}{phosphorylation-deficient receptors still adapt}~\cite{kim1997}.
The phospho-cycle here therefore models only the loss of binding affinity, not the termination of signalling.
Identifying the low-affinity class with the phosphorylated receptor is a simplification, since the affinity conversion (\mk{M1-conv9}{$t_{1/2}\approx9$\,s}~\cite{vanhaastert1984}) is faster than phosphorylation (\mk{M1-phos45}{$t_{1/2}\approx45$\,s}~\cite{vaughan1988}).
Receptor internalisation is not modelled, consistent with the observation that \mk{M1-aff-red}{cAR1 remains at the surface}~\cite{caterina1995b}.

\subsection{Module 2: G-Protein Cycle}\label{sec:m2}

\paragraph{Role.}
cAR1 couples to the heterotrimer G$\alpha_2\beta\gamma$.
FRET between the $\alpha$ and $\beta$ subunits shows that the heterotrimer dissociates and reassociates rapidly on addition and removal of cAMP, that during continuous stimulation the activation reaches a dose-dependent steady state and does \emph{not} decline while the physiological responses subside, and that occupied receptors catalyse the G-protein cycle whether or not they are phosphorylated~\cite{janetopoulos2001}.
Adaptation therefore occurs downstream of the G protein, and free G$\beta\gamma$ is the persistent excitatory signal that Module~3 differentiates.
The initial Ras response to uniform cAMP requires G$\beta$ but not G$\alpha_2$~\cite{kortholt2013}, and symmetry breaking and amplification occur between the G protein and Ras~\cite{kataria2013}.
The module follows the G-protein module of Biswas et al.~\cite{biswas2021}.

\begin{table*}[!t]
\centering
\caption{Module~2 reactions: the G-protein cycle of G$\alpha_2\beta\gamma$. G$\alpha_2^{\mathrm{GTP}}$ and G$\alpha_2^{\mathrm{GDP}}$ are the GTP- and GDP-bound free $\alpha$ subunits; G$\beta\gamma$ is the free dimer. A coloured reaction is the one to which the equally coloured statement in the text and the equally coloured passage in the cited paper refer; a small square appends a further statement.}
\label{tab:m2rxn}
\footnotesize
\renewcommand{\arraystretch}{1.15}
\begin{tabularx}{\textwidth}{@{}l>{\raggedright\arraybackslash}p{0.37\textwidth}>{\raggedright\arraybackslash}X@{}}
\toprule
ID & Reaction & Propensity (rate law)\\
\midrule
\mkt{M2-bis-ke0}{2.0} & \mkt{M2-bis-ke0}{$\mathrm{G}_{\alpha\beta\gamma} \to \mathrm{G}\alpha_2^{\mathrm{GTP}}+\mathrm{G}_{\beta\gamma}$} & $k_{E0}[\mathrm{G}_{\alpha\beta\gamma}]$ \\
\mkn{M2-jan-fret}{M2-bis-g}{2.1} & \mkn{M2-jan-fret}{M2-bis-g}{$\mathrm{G}_{\alpha\beta\gamma}+\mathrm{RC} \to \mathrm{G}\alpha_2^{\mathrm{GTP}}+\mathrm{G}_{\beta\gamma}+\mathrm{RC}$} & $k_E[\mathrm{G}_{\alpha\beta\gamma}][\mathrm{RC}]$ \\
\mkn{M2-jan-phos}{M2-bis-eq}{2.2} & \mkn{M2-jan-phos}{M2-bis-eq}{$\mathrm{G}_{\alpha\beta\gamma}+\mathrm{R}_p\mathrm{C} \to \mathrm{G}\alpha_2^{\mathrm{GTP}}+\mathrm{G}_{\beta\gamma}+\mathrm{R}_p\mathrm{C}$} & $\varepsilon_E k_E[\mathrm{G}_{\alpha\beta\gamma}][\mathrm{R}_p\mathrm{C}]$ \\
\mkt{M2-bis-g}{2.3} & \mkt{M2-bis-g}{$\mathrm{G}\alpha_2^{\mathrm{GTP}} \to \mathrm{G}\alpha_2^{\mathrm{GDP}}$} & $k_\mathrm{hyd}[\mathrm{G}\alpha_2^{\mathrm{GTP}}]$ \\
\mkn{M2-tang-reass}{M2-bis-g}{2.5} & \mkn{M2-tang-reass}{M2-bis-g}{$\mathrm{G}\alpha_2^{\mathrm{GDP}}+\mathrm{G}_{\beta\gamma} \to \mathrm{G}_{\alpha\beta\gamma}$} & $k_\mathrm{re}[\mathrm{G}\alpha_2^{\mathrm{GDP}}][\mathrm{G}_{\beta\gamma}]$ \\
\bottomrule
\end{tabularx}
\end{table*}

\begin{table*}[!t]
\centering
\caption{Module~2 constants.}
\label{tab:m2const}
\footnotesize
\renewcommand{\arraystretch}{1.15}
\begin{tabularx}{\textwidth}{@{}l r l c >{\raggedright\arraybackslash}X@{}}
\toprule
Symbol & Value & Unit & Basis & Meaning, justification, source\\
\midrule
$G{:}R$ & \num{60} &  & E & G-protein to receptor ratio; \mkt{M2-bis-ratio}{Biswas et~al. use 3}~\cite{biswas2021}. Chosen to raise the G$\beta\gamma$ count (shot noise) at constant activated fraction; an explicit trade of abundance realism for signal-to-noise, not a claim about abundance. \\
$G_\mathrm{tot}$ & \num{184620} & $\mathrm{molec/voxel}$ & D & $R_\mathrm{tot}\times G{:}R$ ($38\,\mu$M). \\
$t_{1/2,\mathrm{diss}}$ & \num{3.2} & $\mathrm{s}$ & L & Half-time of FRET loss (G-protein dissociation) after stimulation, \mkt{M2-tang-diss}{$3.4\pm1.4$\,s}~\cite{tang2014}; \mkt{M2-bis-tdiss}{reproduced by the model of}~\cite{biswas2021} (3.0--3.3\,s, Table~2). \\
$k_E$ & \num{0.339} & $\mu\mathrm{M}^{-1}\mathrm{s}^{-1}$ & D & $k_E=\ln 2/(t_{1/2,\mathrm{diss}}\,[\mathrm{R}]_\mathrm{tot})$: pseudo-first-order dissociation when all receptors are occupied. \\
$\varepsilon_E$ & \num{1} &  & L & \mkt{M2-jan-phos}{Phosphorylated and unphosphorylated occupied receptors catalyse the G-protein cycle equally}~\cite{janetopoulos2001}; \mkt{M2-bis-eq}{Biswas et~al. likewise weight all occupied classes equally}~\cite{biswas2021}. \\
$k_{E0}$ & \num{2.18e-7} & $\mathrm{s^{-1}}$ & L & Basal, receptor-independent exchange (\mkt{M2-bis-ke0}{Table~2 of}~\cite{biswas2021}, \mkt{M2-bis-ke0}{value for $G{:}R=3$}~\cite{biswas2021}, not rescaled); sets the resting G$\beta\gamma$ floor. \\
$k_\mathrm{hyd}$ & \num{0.162} & $\mathrm{s^{-1}}$ & C & $7\ln2/30$: intrinsic GTP hydrolysis of G$\alpha_2$, raised 7$\times$ above the value matching the 30\,s recovery anchor to de-saturate the stage (see text). \\
$k_\mathrm{re}$ & \num{6.09} & $\mu\mathrm{M}^{-1}\mathrm{s}^{-1}$ & C & $3500/(t_{1/2,\mathrm{re}}\,[\mathrm{G}]_\mathrm{tot}/2)$ with $t_{1/2,\mathrm{re}}=30$\,s (measured reassociation half-times \mkt{M2-tang-reass}{28.7--32.3\,s}~\cite{tang2014}); the factor 3500 keeps recapture pseudo-first order (see text). The realised reassociation half-time is 7.5\,s, not 30\,s. \\
$D_{\mathrm{G}\alpha\beta\gamma},D_{\mathrm{G}\alpha_2},D_{\mathrm{G}\beta\gamma}$ & \num{0.2} & $\mu\mathrm{m^2/s}$ & E & The value $0.2\,\mu\mathrm{m}^2/\mathrm{s}$ is that of \mkt{M2-bis-D}{Table~3 of}~\cite{biswas2021}; equal mobility of the three forms is a design choice (G proteins also \mkt{M2-elzie}{exchange between cytosol and membrane}~\cite{elzie2009}). \\
\bottomrule
\end{tabularx}
\end{table*}


\paragraph{Reactions.}
\begin{description}\itemsep0pt
\item[2.0] Basal, receptor-independent nucleotide exchange. Its flux is negligible (about 0.03 molecules\,s$^{-1}$ per cell) but it defines the resting G$\beta\gamma$ floor, so that the receiver has a defined baseline rather than a hard zero.
\item[2.1] Receptor-catalysed exchange: an occupied receptor dissociates the heterotrimer into G$\alpha_2$-GTP and G$\beta\gamma$ and is not consumed.
\item[2.2] The same reaction for the phosphorylated occupied receptor. Both forms catalyse the cycle with equal efficiency ($\varepsilon_E=1$), as observed~\cite{janetopoulos2001}, so the phosphorylated receptor carries the drive at high cAMP.
\item[2.3] Intrinsic GTP hydrolysis converts G$\alpha_2$-GTP into G$\alpha_2$-GDP. An RGS-accelerated hydrolysis (2.4) is defined in the code but switched off.
\item[2.5] Recapture of G$\beta\gamma$ by G$\alpha_2$-GDP re-forms the heterotrimer and terminates the signal.
\end{description}

\paragraph{Constants.}
$k_E$ is derived from the dissociation half-time at saturating cAMP: when all receptors are occupied the heterotrimer is consumed with the pseudo-first-order rate $k_E[\mathrm R]_\mathrm{tot}$, so $k_E=\ln2/(t_{1/2,\mathrm{diss}}[\mathrm R]_\mathrm{tot})$ with $t_{1/2,\mathrm{diss}}=3.2$\,s, the half-time of FRET loss measured in living cells, $3.4\pm1.4$\,s~\cite{tang2014}, which the model of Biswas et al.\ reproduces~\cite{biswas2021}.
The basal constant $k_{E0}$ is the value of Table~2 of~\cite{biswas2021} for a ratio of 3.

\emph{Why the gain of the stage was changed.}
The conservation law $[\mathrm G\beta\gamma]=[\mathrm G\alpha_2^\mathrm{GTP}]+[\mathrm G\alpha_2^\mathrm{GDP}]$ holds exactly, because every dissociation creates one G$\alpha_2$ and one G$\beta\gamma$ and every recapture removes one of each.
In the steady state, with $P$ the total exchange flux, $[\mathrm G\alpha_2^\mathrm{GTP}]=P/k_\mathrm{hyd}$ and $k_\mathrm{re}[\mathrm G\alpha_2^\mathrm{GDP}][\mathrm G\beta\gamma]=P$.
If recapture is slow compared with hydrolysis, $[\mathrm G\beta\gamma]\approx\sqrt{P/k_\mathrm{re}}$ and the stage compresses the input contrast (exponent $\tfrac12$); if recapture is fast, $[\mathrm G\beta\gamma]\approx[\mathrm G\alpha_2^\mathrm{GTP}]=P/k_\mathrm{hyd}$ is linear in the drive until the heterotrimer is depleted.
The exponent is therefore a property of the regime and not of the conservation law.
Raising $k_\mathrm{re}$ by a factor of 3500 puts the stage in the linear regime, and raising $k_\mathrm{hyd}$ by a factor of 7 moves the saturation point to higher cAMP.
On the stochastic engine, with replicate cells, the contrast transfer $\mathrm d\ln\mathrm G\beta\gamma/\mathrm d\ln c$ at 60\,nM cAMP increased from 0.22 to 0.77 and the signal-to-noise ratio doubled; at 5\,nM the transfer increased from 0.85 to 1.65 with no loss of signal-to-noise ratio.

\emph{Departures from the literature anchors.}
The chosen values move the two half-time anchors: the dissociation half-time becomes about 2\,s (from 3.2\,s) and the reassociation half-time about 7.5\,s, four times shorter than the measured 29--32\,s~\cite{tang2014}.
The cost of the stochastic simulation at the working point rises about threefold, because the activated fraction is increased on purpose.
The G-protein pool is 60 times the receptor number, 20 times the ratio 3 of~\cite{biswas2021}.
This pool size raises the G$\beta\gamma$ copy number at a fixed activated fraction and was accompanied by paired changes in Module~3 (see there) that leave its dynamics unchanged; it is a deliberate trade of realism for signal-to-noise and is not a claim about G-protein abundance.
The three G-protein forms have the same diffusion constant, which is a design choice.

\subsection{Module 3: Ras Activation, Adaptation and Amplification}\label{sec:m3}

\paragraph{Role.}
The earliest polarised response downstream of the G protein is the activation of Ras.
RasG and RasC are \mk{M3-trans}{activated rapidly and transiently by cAMP in aggregation-competent cells}~\cite{kae2004}, and \mk{M3-pi3k}{RasG is the upstream activator of PI3K}~\cite{sasaki2004,takeda2012}.
Adaptation takes place between the G protein and Ras: \mk{M3-g-persist}{G-protein activation does not decline during continuous stimulation, whereas the downstream responses subside}~\cite{janetopoulos2001,elzie2009}, and \mk{M3-g-steps}{under two successive cAMP steps the G-protein dissociation rises step-like and persists, while Ras responds transiently to each step}~\cite{xu2022}.
\mk{M3-adapt}{The activation of RasG shows near-perfect adaptation over a wide range of concentrations}~\cite{takeda2012}; among the two simple topologies that can produce perfect adaptation, \mk{M3-iffl}{only an incoherent feed-forward loop fitted the measured responses}~\cite{takeda2012}, in which \mk{M3-prop}{the activator drives an excitatory and an inhibitory node proportionally}~\cite{goentoro2009,takeda2012}, \mk{M3-gap-glob}{possibly with a RasGAP as the global inhibitor}~\cite{takeda2012}.
An incoherent feed-forward loop of this type has the structure of a local-excitation--global-inhibition (LEGI) gradient sensor~\cite{levchenko2002,ma2004}.
In the Ras network, \mk{M3-legi}{the activator and the inactivator of Ras, a RasGEF and a RasGAP, are both activated by receptor and G-protein signalling}~\cite{takeda2012,meierschellersheim2006}, and such a loop can provide \mk{M3-fcd3}{fold-change detection}~\cite{goentoro2009}.
The RasGAP NF1 (DdNF1) is \mk{M3-nf1}{a major regulator of Ras activity and its loss makes Ras activity spatially and temporally unregulated}~\cite{zhang2008}; \mk{M3-nf1-rasg}{in its absence the activation of RasG is delayed and extended, whereas RasD, Rap1 and RasC are unaffected}~\cite{zhang2008}.
\mk{M3-rect}{Ras activation is suppressed when the cAMP concentration falls (``rectification''), which arises naturally in an incoherent feed-forward loop with zero-order ultrasensitivity}~\cite{nakajima2014}.
\mk{M3-ras-no-down}{The Ras response to uniform cAMP needs none of the downstream PI3K, TorC2, PLA2 and sGC pathways and is unchanged when all four are inhibited, also together with latrunculin~A and in a pkbR1-null background}~\cite{kortholt2011}.
Finally, the Ras/PI3K/PIP3 network is \mk{M3-excit}{excitable: it fires without external cues}~\cite{devreotes2017,fukushima2019,miao2017}, \mk{M3-bias}{a stimulus biases threshold and timing}~\cite{fukushima2019,miao2017,devreotes2017,iglesias2012}, and \mk{M3-excit-intr}{Ras remains excitable without the PIP3, TorC2, PLA2 and sGC pathways and without a functional actin cytoskeleton}~\cite{fukushima2019,iwamoto2025}.
\mk{M3-dnf}{Excitability requires a positive feedback for the all-or-none response and a delayed negative feedback that terminates the response and is followed by a refractory period}~\cite{fukushima2019}.
Module~3 therefore combines (i)~a fast G$\beta\gamma$-driven excitor and a slower G$\beta\gamma$-driven inhibitor (adaptation and spatial comparison), (ii)~a saturable GAP (rectification), (iii)~a positive feedback from PIP3 onto RasG (amplification) and (iv)~a delayed negative feedback read from RasG-GTP, the brake (termination of excursions that the stimulus did not cause; proposed).
An integral-feedback (antithetic) arm that was tested earlier was removed, because \mk{M3-noint}{integral feedback gives oscillations and stimulus-dependent kinetics that are inconsistent with the measured Ras response}~\cite{takeda2012}.

\paragraph{Adaptation and brake.}
The two negative arms act on the same RasG-GTP but read different variables, and this is what separates them.
The inhibitor of the incoherent feed-forward loop (3.4e, 3.5, 3.7z) reads the input, G$\beta\gamma$, as the excitor does.
It sets the adapted level, which is therefore independent of the dose, and it alone acts when the stimulus falls: \mk{M3-undershoot}{after a sudden decrease of cAMP the Ras reporter returns rapidly to the cytosol and then slowly to its basal amount}~\cite{takeda2012}, i.e.\ RasG-GTP transiently falls below its resting level because the slow inhibitor outlasts the fast excitor.
The brake (3.3c$_\mathrm{R}$, 3.3d, 3.3e) reads the output, RasG-GTP.
Its purpose is the delayed negative feedback that excitability requires (see Role): it is the only arm that can terminate a Ras excursion that the stimulus did not cause, such as a spontaneous firing of the excitable network, because the inhibitor follows G$\beta\gamma$ and is not raised by such an excursion.
After a brief stimulus both arms contribute to the refractory period, since B$^*$ and the inhibitor both decay with $\tau=10$\,s; the measured target is that \mk{M3-refr}{with paired 2-s mechanical stimuli and an F-actin reporter, the second stimulus elicits no response at intervals below 12\,s, responsiveness then recovers with a half-time of about 7\,s as observed for chemoattractant, and a refractory period is one of the features of excitable systems}~\cite{artemenko2016}.
The brake is constructed so that it does not set the adapted level (design choice): B$^*$ decays in first order and has no set point, and its half-activation lies at the driven RasG-GTP level, so that its removal flux, which is proportional to $[\mathrm{B}^*][\mathrm{RasG^{GTP}}]$ and hence roughly to the square of RasG-GTP while B is far from saturation, is small at rest and in the adapted plateau and large only at the peak.
It is not the integral controller that was rejected for Ras: in that model \mk{M3-int-def}{the GAP is produced in proportion to Ras-GTP and removed at a constant, zero-order rate}~\cite[Suppl.\ Information]{takeda2012}, which holds Ras-GTP at a set point, whereas B$^*$ is removed in first order.
That the brake does not adapt is a deviation from the literature, in which a feedback read from Ras is proposed as an adaptation mechanism: \mk{M3-nfblb-def}{a negative feedback loop with a buffer node, in which the output is shut down by an inhibitor induced by the output itself, is one of the two simple topologies that produce adaptation}~\cite{xu2022}, and \mk{M3-nfblb-fit}{models in which the RasGAP is activated by G$\alpha_2$-GTP, by Ras-GTP or by both reproduce transient, adaptive Ras signalling}~\cite{xu2022}.
The model assigns adaptation to the incoherent feed-forward loop, the only topology that fitted the measured responses in the comparison with integral feedback (see Role), and keeps the Ras-read feedback out of the adapted state; this division of labour is a design choice and is not shown by the sources.
The absence of a second excursion after a step is a test of the tuning.
During a sustained cAMP step both arms act; they are distinguished by the undershoot after removal of the stimulus (inhibitor only) and by the termination of spontaneous excursions (brake only).

\begin{table*}[!t]
\centering
\caption{Module~3 reactions in the default configuration. GEF$_R$ and GAP are the cytosolic RasG exchange factor and GTPase-activating protein (asterisk: active form); $\mathrm{RasG}^{\mathrm{GTP/GDP}}$ are the membrane-bound Ras states; B is the substrate of the delayed brake (3.3c$_\mathrm{R}$--e), activated by RasG-GTP. 3.1--3.7z form the adaptation (incoherent feed-forward loop, read from G$\beta\gamma$), 3.3b the amplification and 3.3c$_\mathrm{R}$--e the brake (read from RasG-GTP). Reactions marked $^\dagger$ are \emph{proposed}: the mechanism is a candidate, and they are kept only if tuning shows that they are needed (see text). A coloured reaction is the one to which the equally coloured statement in the text and the equally coloured passage in the cited paper refer; a small square appends a further statement.}
\label{tab:m3rxn}
\footnotesize
\renewcommand{\arraystretch}{1.15}
\begin{tabularx}{\textwidth}{@{}l>{\raggedright\arraybackslash}p{0.37\textwidth}>{\raggedright\arraybackslash}X@{}}
\toprule
ID & Reaction & Propensity (rate law)\\
\midrule
\mkn{M3-kort}{M3-prop,M3-legi}{3.1} & \mkn{M3-kort}{M3-prop,M3-legi}{$\mathrm{G}_{\beta\gamma}+\mathrm{GEF_R} \to \mathrm{GEF_R^*}+\mathrm{G}_{\beta\gamma}$} & $k_{E,\mathrm{on}}[\mathrm{G}_{\beta\gamma}][\mathrm{GEF_R}]$ \\
\mkn{M3-peak-fast}{M3-tau-fit}{3.2} & \mkn{M3-peak-fast}{M3-tau-fit}{$\mathrm{GEF_R^*} \to \mathrm{GEF_R}$} & $k_{E,\mathrm{off}}[\mathrm{GEF_R^*}]$ \\
\mkn{M3-gefr}{M3-trans,M3-gef-exch}{3.3} & \mkn{M3-gefr}{M3-trans,M3-gef-exch}{$\mathrm{GEF_R^*}+\mathrm{RasG}^{\mathrm{GDP}} \to \mathrm{RasG}^{\mathrm{GTP}}+\mathrm{GEF_R^*}$} & $k_{R,\mathrm{on}}[\mathrm{GEF_R^*}][\mathrm{RasG}^{\mathrm{GDP}}]$ \\
\mkn{M3-prop}{M3-gap-glob,M3-legi,M3-iffl,M3-gap-gbg}{3.4e} & \mkn{M3-prop}{M3-gap-glob,M3-legi,M3-iffl,M3-gap-gbg}{$\mathrm{G}_{\beta\gamma}+\mathrm{GAP} \to \mathrm{GAP^*}+\mathrm{G}_{\beta\gamma}$} & $k_{I,\mathrm{on}}[\mathrm{G}_{\beta\gamma}][\mathrm{GAP}]$ \\
\mkt{M3-basal-gap}{3.4b} & \mkt{M3-basal-gap}{$\mathrm{GAP} \to \mathrm{GAP^*}$} & $k_{I,\mathrm{b}}[\mathrm{GAP}]$ \\
\mkn{M3-tau-fit}{M3-slower-inh,M3-undershoot}{3.5} & \mkn{M3-tau-fit}{M3-slower-inh,M3-undershoot}{$\mathrm{GAP^*} \to \mathrm{GAP}$} & $k_{I,\mathrm{off}}[\mathrm{GAP^*}]$ \\
\mkn{M3-rect}{M3-nf1kin,M3-nf1-rasg}{3.7z} & \mkn{M3-rect}{M3-nf1kin,M3-nf1-rasg}{$\mathrm{GAP^*}+\mathrm{RasG}^{\mathrm{GTP}} \to \mathrm{RasG}^{\mathrm{GDP}}+\mathrm{GAP^*}$} & $k_\mathrm{cat}[\mathrm{GAP^*}]\,\dfrac{[\mathrm{RasG}^{\mathrm{GTP}}]}{K_\mathrm{NF1}+[\mathrm{RasG}^{\mathrm{GTP}}]}$ \\
\mkn{M3-pfb}{M3-pip3-gef,M3-ly-red,M3-pip3-small}{3.3b} & \mkn{M3-pfb}{M3-pip3-gef,M3-ly-red,M3-pip3-small}{$\mathrm{RasG}^{\mathrm{GDP}}+\mathrm{PIP3} \to \mathrm{RasG}^{\mathrm{GTP}}+\mathrm{PIP3}$} & $k_\mathrm{PIP3,R}[\mathrm{PIP3}][\mathrm{RasG}^{\mathrm{GDP}}]$ \\
\mkn{M3-nfblb-def}{M3-c2gap-ras,M3-dnf}{3.3c$_\mathrm{R}^\dagger$} & \mkn{M3-nfblb-def}{M3-c2gap-ras,M3-dnf}{$\mathrm{RasG}^{\mathrm{GTP}}+\mathrm{B} \to \mathrm{B}^*+\mathrm{RasG}^{\mathrm{GTP}}$} & $k_{B,\mathrm{on}}[\mathrm{RasG}^{\mathrm{GTP}}][\mathrm{B}]$ \\
\mkn{M3-refr}{M3-slow-buf}{3.3d$^\dagger$} & \mkn{M3-refr}{M3-slow-buf}{$\mathrm{B}^* \to \mathrm{B}$} & $k_{B,\mathrm{off}}[\mathrm{B}^*]$ \\
\mkn{M3-c2gap-ras}{M3-nfblb-fit}{3.3e$^\dagger$} & \mkn{M3-c2gap-ras}{M3-nfblb-fit}{$\mathrm{B}^*+\mathrm{RasG}^{\mathrm{GTP}} \to \mathrm{RasG}^{\mathrm{GDP}}+\mathrm{B}^*$} & $k_{B,\mathrm{hyd}}[\mathrm{B}^*][\mathrm{RasG}^{\mathrm{GTP}}]$ \\
\bottomrule
\end{tabularx}
\end{table*}

\begin{table*}[!t]
\centering
\caption{Module~3 constants. Rate constants and pool sizes are paired: $k_{E,\mathrm{on}}[\mathrm{GEF_R}]_\mathrm{tot}$, $k_{R,\mathrm{on}}[\mathrm{GEF_R}]_\mathrm{tot}$ and $k_{I,\mathrm{on}}[\mathrm{GAP}]_\mathrm{tot}$ were held fixed when the pools were enlarged, so that only shot noise changed.}
\label{tab:m3const}
\footnotesize
\renewcommand{\arraystretch}{1.15}
\begin{tabularx}{\textwidth}{@{}l r l c >{\raggedright\arraybackslash}X@{}}
\toprule
Symbol & Value & Unit & Basis & Meaning, justification, source\\
\midrule
$[\mathrm{RasG}]_\mathrm{tot}$ & \num{38500} & $\mathrm{molec/voxel}$ & E & $5\times10^5$ per cell; large for low shot noise. The copy number is not well measured. \\
$[\mathrm{GEF_R}]_\mathrm{tot}$ & \num{138600} & $\mathrm{molec/voxel}$ & E & Excitor pool ($1.8\times10^6$ per cell). \mkt{M3-gefr}{RasGEFR activates RasG}~\cite{kae2007}. \\
$[\mathrm{GAP}]_\mathrm{tot}$ & \num{46200} & $\mathrm{molec/voxel}$ & E & Inhibitor pool ($6\times10^5$ per cell); \mkt{M3-nf1}{a RasGAP (NF1) is the global inhibitor}~\cite{zhang2008}. \\
$k_{E,\mathrm{on}}$ & \num{3.8e-3} & $\mu\mathrm{M}^{-1}\mathrm{s}^{-1}$ & C & G$\beta\gamma$-driven GEF activation; \mkt{M3-kort}{the initial Ras response requires G$\beta$}~\cite{kortholt2013}. \\
$k_{E,\mathrm{off}}$ & \num{0.5} & $\mathrm{s^{-1}}$ & C & GEF inactivation, $\tau_E=2$\,s: \mkt{M3-peak-fast}{the excitor is fast, so the Ras peak occurs within seconds}~\cite{kae2004,takeda2012}; close to \mkt{M3-tau-fit}{the fitted $0.4\,\mathrm{s^{-1}}$}~\cite{takeda2012}. \\
$k_{R,\mathrm{on}}$ & \num{0.333} & $\mu\mathrm{M}^{-1}\mathrm{s}^{-1}$ & C & GEF-catalysed GDP$\to$GTP exchange on RasG. \\
$k_{I,\mathrm{on}}$ & \num{7.67e-4} & $\mu\mathrm{M}^{-1}\mathrm{s}^{-1}$ & D & G$\beta\gamma$-driven GAP activation. It reads the \emph{same} input as 3.1, as \mkt{M3-prop}{an incoherent feed-forward loop requires}~\cite{takeda2012,goentoro2009}. Calibrated value $2.53\times10^{-4}$ multiplied by $3.03$ together with $k_{I,\mathrm{off}}$, which keeps the activated fraction; \mkt{M3-tau-pair}{the fit tied activation and deactivation rates in the same way}~\cite{takeda2012}. \\
$k_{I,\mathrm{off}}$ & \num{0.1} & $\mathrm{s^{-1}}$ & L & GAP inactivation, $\tau_I=10$\,s, \mkt{M3-tau-fit}{the fitted value, which sets the time of the return to the basal amount}~\cite{takeda2012} (until 2026-10: $0.033\,\mathrm{s^{-1}}$); \mkt{M3-slower-inh}{the inhibitor is slower than the excitor, which produces adaptation}~\cite{takeda2012}. Not yet recalibrated against the downstream modules. \\
$k_{I,\mathrm{b}}$ & \num{1.52e-3} & $\mathrm{s^{-1}}$ & D & Stimulus-independent GAP activity. Without it the resting Ras/PIP3 loop latches; consistent with \mkt{M3-basal-gap}{a basal RasGAP}~\cite{zhang2008,xu2021c2gap}. Calibrated as the minimum value giving the low rest state ($5\times10^{-4}$ at $\tau_I=30$\,s), multiplied by $3.03$ with $k_{I,\mathrm{off}}$, which keeps the basal activated fraction. \\
$K_\mathrm{NF1}$ & \num{0.2} & $\mu\mathrm{M}$ & L & \mkt{M3-nf1kin}{Michaelis constant of the NF1 GAP domain, $0.3\,\mu$M}~\cite{wiesmuller1992,coyle2016}; the value is rounded to $0.2$. Mammalian catalytic domain; no \emph{Dictyostelium} kinetics exist. \\
$k_\mathrm{cat}$ & \num{15.8} & $\mathrm{s^{-1}}$ & D & Derived from the resting distance to the zero-order switch, $\theta_\mathrm{rest}=0.5$. Only the product $k_\mathrm{cat}[\mathrm{GAP}^*]$ enters; it is 11$\times$ the mammalian \mkt{M3-nf1kin}{$k_\mathrm{cat}=1.4$\,s$^{-1}$}~\cite{coyle2016,wiesmuller1992}, so the active-GAP pool stands for more enzyme than is measured. \\
$k_\mathrm{PIP3,R}$ & \num{0.04} & $\mu\mathrm{M}^{-1}\mathrm{s}^{-1}$ & C & PIP3$\to$RasG feedback (3.3b), \mkt{M3-pip3-gef}{the PIP3 term of the Ras GEF reaction}~\cite{fukushima2019}. The value committed for this form in an earlier calibration (higher gains raised the variance but not the contrast, see text); not re-tuned after the changes of 2026-10. Kept small, since \mkt{M3-ly-red}{PI3K inhibition reduces but does not abolish Ras activation}~\cite{sasaki2004}. Expressed relative to the legacy PIP2 pool (Sec.~\ref{sec:m4}). \\
$k_{B,\mathrm{on}}$ & \num{0.0753} & $\mu\mathrm{M}^{-1}\mathrm{s}^{-1}$ & E & Brake activation by RasG-GTP (3.3c$_\mathrm{R}$); half-activation at 6400 RasG-GTP molecules/voxel, the driven level recorded in the code (design anchor, not tuned). Proposed brake. \\
$k_{B,\mathrm{off}}$ & \num{0.1} & $\mathrm{s^{-1}}$ & D & Brake decay ($\tau=10$\,s), from \mkt{M3-refr}{the recovery half-time of about 7\,s after a refractory period}~\cite{artemenko2016} as $\tau=t_{1/2}/\ln2$, assuming that B$^*$ governs the recovery; measured with an F-actin reporter after mechanical stimulation, not for Ras. \\
$k_{B,\mathrm{hyd}}$ & \num{39.8} & $\mu\mathrm{M}^{-1}\mathrm{s}^{-1}$ & D & Chosen so that at half activation the brake removes RasG-GTP at $0.6\times$ the rate of the LEGI GAP arm at the driven state (design target, not tuned). Proposed brake. \\
$[\mathrm{B}]_\mathrm{tot}$ & \num{1540} & $\mathrm{molec/voxel}$ & E & Brake pool; the mean braking action is independent of pool size, so it was raised to reduce noise. \\
$D_{\mathrm{GEF_R}},D_{\mathrm{RasG}}$ & \num{0.1} & $\mu\mathrm{m^2/s}$ & E & Local excitor and membrane-bound Ras (prenylated anchor). \\
$D_{\mathrm{GAP}},D_\mathrm{B}$ & \num{20} & $\mu\mathrm{m^2/s}$ & E & Global inhibitor and global brake. Cytosolic value; for B$^*$, a membrane-attached species in the code, it contradicts the local purpose of the brake (open, see 3.3c$_\mathrm{R}$--e); $0.4\,\mu\mathrm{m^2/s}$ gives the local brake. \\
\bottomrule
\end{tabularx}
\end{table*}


\paragraph{Reactions.}
\begin{description}\itemsep0pt
\item[3.1/3.2] The excitor GEF$_\mathrm{R}$ is activated by free G$\beta\gamma$ and inactivated with $\tau_E=2$\,s. \mk{M3-gefr}{RasGEFR is the exchange factor of RasG}~\cite{kae2007}, and \mk{M3-kort}{the initial Ras response depends on G$\beta$}~\cite{kortholt2013}. The short relaxation time makes the RasG peak follow the input within seconds. It is close to the fitted value: \mk{M3-tau-fit}{in the fitted incoherent feed-forward model the GEF is inactivated at $0.4\,\mathrm{s^{-1}}$ and the GAP at $0.1\,\mathrm{s^{-1}}$, the GAP rate sets the time of the return to the basal amount, and the GEF rate must exceed the GAP rate}~\cite[Suppl.\ Information]{takeda2012}.
\item[3.3] The active GEF exchanges GDP for GTP on membrane-bound RasG, which produces RasG-GTP. This is the general mechanism of Ras activation: \mk{M3-gef-exch}{Ras GEFs catalyse the exchange of GDP for GTP and thereby activate Ras}~\cite{kae2004,kae2007,xu2022}.
\item[3.4e/3.4b/3.5] The inhibitor GAP is activated by G$\beta\gamma$ (3.4e), i.e.\ by the same variable as the excitor, and decays with $\tau_I=10$\,s (3.5), the GAP rate of the fit quoted under 3.1/3.2 (until 2026-10 the model used $\tau_I=30$\,s). A G$\beta\gamma$-armed GAP is also the choice of an earlier model, in which \mk{M3-gap-gbg}{a RasGAP is activated by G$\beta\gamma$, translocates to the membrane and deactivates Ras}~\cite{meierschellersheim2006}. The requirement that both arms read the same input is essential: if they read different upstream variables with different dose dependence, the ratio of excitation to inhibition acquires a dose dependence of its own and the stage inverts instead of adapting (this was found when the inhibitor read the receptor occupancy). A small basal activation (3.4b) represents the stimulus-independent GAP at rest, which is needed because the resting Ras/PIP3 loop otherwise latches on; a basal RasGAP is consistent with \mk{M3-basal-gap}{the enhanced basal Ras activity of GAP-deficient cells}~\cite{zhang2008,xu2021c2gap}.
\item[3.7z] The active GAP hydrolyses RasG-GTP with Michaelis--Menten kinetics. Because $K_\mathrm{NF1}$ is far below the Ras pool, the enzyme operates in the zero-order regime, where a graded input is turned into a switch-like output and \mk{M3-rect}{the response to a falling stimulus is suppressed}~\cite{nakajima2014}. \mk{M3-nf1kin}{The Michaelis constant of the NF1 GAP domain is $0.3\,\mu$M and its $k_\mathrm{cat}$ is $1.4$\,s$^{-1}$ (mammalian catalytic domain)}~\cite{wiesmuller1992,coyle2016}. The plain mass-action form of this step is available and is not used.
\item[3.3b] PIP3 converts membrane-bound RasG-GDP into RasG-GTP. Together with 4.1 (RasG-GTP recruits PI3K) and 4.3 (PI3K makes PIP3) this closes the \mk{M3-pfb}{Ras$\to$PI3K$\to$PIP3$\to$Ras positive-feedback loop, which is the amplifier of the excitable network}~\cite{fukushima2019}, consistent with \mk{M3-kat}{the amplification between the G protein and Ras}~\cite{kataria2013}. The form is that of \mk{M3-pip3-gef}{the published Ras/PIP3 model, whose Ras GEF reaction has a basal term and a term for the feedback from PIP3}~\cite{fukushima2019}. The feedback contributes but is not required: \mk{M3-ly-red}{initial Ras activation requires neither PI3K activity nor F-actin, but its level is reduced by LY294002 or latrunculin}~\cite{sasaki2004}, \mk{M3-ras-pip3-indep}{PIP3 production and degradation are not necessary for Ras wave generation}~\cite{fukushima2019}, and \mk{M3-pip3-small}{elevated Ras activity depends largely on decreased PI(3,4)P$_2$, with a contribution of feedback from PIP3}~\cite{li2018}.
The molecule on which PIP3 acts is not identified; a route through GAPs is possible, since \mk{M3-li-gap}{RasGAP2 and RapGAP3 bind PI(3,4)P$_2$}~\cite{li2018}, and the direct exchange of 3.3b is the lumped form.
The feedback is not placed on the G protein, i.e.\ as a PIP3-accelerated dissociation of the heterotrimer (3.3g, available as an option), because \mk{M3-g-follows-stim}{G-protein activation provides a simple intracellular translation of the external gradient and reflects the local cAMP receptor occupancy}~\cite{xu2005} and not PIP3 (see Role), \mk{M3-gbg-pi3k-indep}{the cAMP-induced change of the G$\beta\gamma$ mobility and domain formation are independent of PI3K activity}~\cite{vanhemert2010} (mobility, not dissociation, was measured), and \mk{M3-sb-pip3-indep}{symmetry breaking of Ras requires G$\alpha_2$ and G$\beta\gamma$ but not the PIP3, cGMP, TorC2 and PLA2 pathways}~\cite{kortholt2013}.
A second positive feedback may act within the Ras network itself: \mk{M3-gef-pfb}{the co-localisation of RasGEFX and RasGEFB with Ras-GTP-enriched domains suggests that their recruitment depends on Ras-GTP, and hence a positive feedback between these GEFs and Ras-GTP}~\cite{iwamoto2025}. It is not modelled, so the excitability of the model runs through 3.3b only, whereas the cell is excitable also without PIP3 (see Role).
\item[3.3c$_\mathrm{R}$--e] \textbf{Brake: delayed negative feedback read from RasG-GTP; the mechanism is a candidate, the molecule is not identified.}
RasG-GTP activates the brake substrate B (3.3c$_\mathrm{R}$), B$^*$ decays with $\tau=10$\,s (3.3d) and accelerates the hydrolysis of RasG-GTP (3.3e).
\emph{Purpose.} The brake supplies the delayed negative feedback that excitability requires and terminates excursions that the stimulus did not cause (see Role and Adaptation and brake).
\emph{Reader.} The feedback is read from RasG-GTP and not from PIP3 or a PKB, because Ras remains excitable, and its response to cAMP is unchanged, without the PIP3, TorC2, PLA2 and sGC pathways (see Role); a feedback that terminates a Ras excursion therefore cannot require them.
Placing the feedback on Ras follows the published model, but not its form: \mk{M3-fuk-nfb}{that model places the positive and the delayed negative feedback on Ras, and its GAP recruitment to the membrane is promoted by Ras-GDP and opposed by Ras-GTP}~\cite{fukushima2019}; a GAP activity induced by RasG-GTP, as in 3.3c$_\mathrm{R}$, is the simpler form chosen here.
\emph{Candidate molecule.} A RasGAP that is activated through Ras exists: \mk{M3-c2gap-ras}{membrane translocation and activation of the RasGAP C2GAP1 require Ras on the membrane, i.e.\ a negative feedback with a buffer node}~\cite{xu2022} (the primary study is cited there and is not among the sources used here).
In the sources, however, it has a different role from the brake: \mk{M3-c2gap-adapt}{C2GAP1 controls both the basal Ras activity and the GPCR-mediated Ras adaptation}~\cite{xu2021c2gap}, whereas the brake is off at rest and does not set the adapted level, and whether C2GAP1 acts on RasG in particular is not shown.
C2GAP1 is therefore a candidate for the molecule, not evidence for the role; that a Ras-read GAP terminates spontaneous excursions is an assumption of the model.
\emph{Time constant.} $\tau$ is set by the refractory period: a recovery half-time of about 7\,s (see Adaptation and brake) corresponds to $\tau=t_{1/2}/\ln 2\approx10$\,s, assuming that the recovery is governed by the decay of B$^*$.
The half-time was measured with an F-actin reporter after mechanical stimulation, not for Ras, and since $\tau_I$ has the same value, the measurement does not separate the brake from the inhibitor.
\emph{Range (open contradiction).} The purpose and the candidate are local: a spontaneous excursion is a local event, and \mk{M3-c2gap-local}{C2GAP1 is a locally recruited RasGAP}~\cite{xu2021c2gap}. In the model, however, B and B$^*$ have the cytosolic diffusion constant of $20\,\mu\mathrm{m^2/s}$, so that within $\tau=10$\,s B$^*$ spreads over $\sqrt{D\tau}\approx14\,\mu$m, about one cell length, and the brake acts globally.
The global value was set for the earlier PIP3 reader, after a local brake ($D=0.4\,\mu\mathrm{m^2/s}$) had reduced the information about the gradient direction at the Ras stage to less than half of that without a brake; a local brake is strongest where the signal is strongest and thereby cancels the contrast, whereas a global brake subtracts the same amount everywhere (argued, not measured).
A local brake would terminate an excursion where it occurs but, by that comparison, costs gradient contrast; whether this also holds for the RasG reader has not been tested. The range is therefore a design choice that tuning has to decide (\texttt{DICTY\_M3\_DBRAKE}).
\emph{Alternative: PKB.} \mk{M3-miao-ext}{The negative regulation by the PKBs extends to other Ras proteins and to PI3K, which makes them candidates for the delayed negative feedback of the excitable network}~\cite{miao2017}, but \mk{M3-miao-mech}{they are suggested to act indirectly, by inhibiting the RasGEF Aimless and by activating PI5K}~\cite{miao2017}, and the documented PKB feedback spares RasG: \mk{M3-sca1-rasc}{the Sca1 complex regulates RasC activity while RasG activation is unaffected}~\cite{charest2010}. That edge is modelled in Module~8 (8.2p/8.2q/8.2s). The PKB reader (3.3c$'$) and the PIP3 reader (3.3c) remain available as options.
\emph{Status.} The brake is a proposed term. Its half-activation and strength are design targets that have not yet been tuned; tuning has to show that the rest state is the low branch, that the adapted plateau is not shifted by the brake, that no second excursion follows a step, and that the module resets as measured: \mk{M3-reset}{after removal of cAMP the network returns to its prestimulus state in less than 1\,min}~\cite{xu2022}.
\end{description}

\paragraph{Constants.}
The rate constants of the LEGI arms (3.1--3.5) are constrained by dynamical targets rather than by single measurements: an excitor time constant of 2\,s, an inhibitor time constant of 10\,s, both from the fit of the measured Ras response, and near-perfect adaptation of the plateau to a sustained step.
When $\tau_I$ was shortened from 30\,s to 10\,s, $k_{I,\mathrm{on}}$ and $k_{I,\mathrm{b}}$ were multiplied by the same factor $0.1/0.033=3.03$ as $k_{I,\mathrm{off}}$, so that all activated GAP fractions, and with them $k_\mathrm{cat}$, are unchanged and only the time constant moves; the fit tied the rates in the same way, as \mk{M3-tau-pair}{the activation rates were set to one tenth of the deactivation rates}~\cite[Suppl.\ Information]{takeda2012}. This change has not yet been recalibrated against the downstream modules.
The pool sizes were increased to reduce shot noise, and the rate constants were rescaled so that only the products $k_{E,\mathrm{on}}[\mathrm{GEF_R}]_\mathrm{tot}$, $k_{R,\mathrm{on}}[\mathrm{GEF_R}]_\mathrm{tot}$ and $k_{I,\mathrm{on}}[\mathrm{GAP}]_\mathrm{tot}$ enter the dynamics; the pool enlargement in Module~2 was paired with rescaling of $k_{E,\mathrm{on}}$ and $k_{I,\mathrm{on}}$ so that the activated fractions and time constants do not change.
Unpaired, the enlargement makes the stage worse, since the activated fraction of the excitor changes with the G$\beta\gamma$ level.
$K_\mathrm{NF1}$ is taken from the mammalian catalytic domain because no \emph{Dictyostelium} kinetics exist; the $k_\mathrm{cat}$ of the model is derived from the distance to the zero-order switch and is 11 times the mammalian value, so the active-GAP pool stands for more enzyme than is measured.
The feedback gain $k_\mathrm{PIP3,R}$ cannot be raised as a means to improve gradient contrast: because the feedback acts on the level, it lifts the resting baseline as much as the driven response, and a contrast is a ratio; above a critical gain the loop is bistable and stops resetting.
Of the brake constants only the decay has an anchor (the refractory period); the activation is derived from a half-activation at 6400 RasG-GTP molecules per voxel, the driven level recorded in the code, and the removal rate from a target strength (0.6 of the LEGI GAP arm at the driven state), so that the brake is off at rest and engages at the peak.
The brake and the GAP have a cytosolic diffusion constant of $20\,\mu\mathrm{m^2/s}$, which is appropriate for a cytosolic GAP but not for the membrane-associated brake B$^*$ of the code; it makes the brake global rather than local, which contradicts its local purpose and its local candidate molecule (see 3.3c$_\mathrm{R}$--e, Range) and is an open point of the present parameterisation.
The saturable GAP is discussed in Table~\ref{tab:hill}.

\subsection{Module 4: PIP3, PTEN and the PkbA Feedback}\label{sec:m4}

\paragraph{Role.}
In chemotaxing cells \mk{M4-recip}{PI3K moves to the plasma membrane at the front and PTEN shows a reciprocal distribution at the back}~\cite{funamoto2002,iijima2002}.
\mk{M4-pip3-trans}{Bulk PIP3 rises transiently after a stimulus, PI3K activation increases within 5\,s and then declines}, and \mk{M4-pten-dephos}{the changes are larger and longer in cells lacking PTEN}~\cite{huang2003}.
PIP3 is read out by PH-domain proteins such as CRAC, \mk{M4-crac}{whose translocation reflects activation of the G-protein pathway}~\cite{parent1998}; after uniform stimulation \mk{M4-phcrac}{PHCrac-GFP translocation peaks after 6--8\,s and disappears after about 20\,s, followed by small patches from about 30\,s}~\cite{postma2003}.
\mk{M4-ras-pi3k}{Ras activates PI3K}~\cite{sasaki2004,fukushima2019}, and \mk{M4-mutual}{PIP3 in turn suppresses the membrane binding of PTEN, so that PIP3 and PTEN are mutually inhibitory and can form a bistable pair}~\cite{matsuoka2018}.
\mk{M4-pkb-delay}{A delayed negative feedback that reduces Ras and PI3K activity is provided by the PKBs, whose activation lags behind that of Ras and PI3K}~\cite{miao2017,charest2010}.
Module~4 implements these elements with a conserved lipid pool, a conserved PTEN pool and a two-step PkbA arm.

\begin{table*}[!t]
\centering
\caption{Module~4 reactions: PI3K recruitment, PIP3/PTEN mutual inhibition, PIP3 removal, and the delayed PkbA feedback. Subscripts c and m denote cytosolic and membrane-bound forms; $N$ is a molecule count per voxel.}
\label{tab:m4rxn}
\footnotesize
\renewcommand{\arraystretch}{1.15}
\begin{tabularx}{\textwidth}{@{}l>{\raggedright\arraybackslash}p{0.37\textwidth}>{\raggedright\arraybackslash}X@{}}
\toprule
ID & Reaction & Propensity (rate law)\\
\midrule
4.1 & $\mathrm{RasG}^{\mathrm{GTP}}+\mathrm{PI3K}_c \to \mathrm{PI3K}_m+\mathrm{RasG}^{\mathrm{GTP}}$ & $k_{3K,\mathrm{on}}[\mathrm{RasG}^{\mathrm{GTP}}][\mathrm{PI3K}_c]$ \\
4.2 & $\mathrm{PI3K}_m \to \mathrm{PI3K}_c$ & $k_{3K,\mathrm{off}}[\mathrm{PI3K}_m]$ \\
4.3 & $\mathrm{PIP2}+\mathrm{PI3K}_m \to \mathrm{PIP3}+\mathrm{PI3K}_m$ & $k_\mathrm{syn}[\mathrm{PIP2}][\mathrm{PI3K}_m]$ \\
4.4 & $\mathrm{PTEN}_c+\mathrm{PIP2} \to \mathrm{PTEN}_m+\mathrm{PIP2}$ & $k_{PT,\mathrm{on}}[\mathrm{PTEN}_c][\mathrm{PIP2}]$ \\
4.5 & $\mathrm{PTEN}_m \to \mathrm{PTEN}_c$ & $k_{PT,\mathrm{off}}[\mathrm{PTEN}_m]$ \\
4.6c & $\mathrm{PTEN}_m \to \mathrm{PTEN}_c$ & $k_\mathrm{ev}[\mathrm{PTEN}_m]\,\dfrac{N_\mathrm{PIP3}^2}{K_\mathrm{ev}^2+N_\mathrm{PIP3}^2}$ \\
4.7 & $\mathrm{PIP3}+\mathrm{PTEN}_m \to \mathrm{PIP2}+\mathrm{PTEN}_m$ & $k_{PT}[\mathrm{PIP3}][\mathrm{PTEN}_m]$ \\
4.8 & $\mathrm{PIP3} \to \mathrm{PIP2}$ & $k_\mathrm{SH}[\mathrm{PIP3}]$ \\
4.10 & $\mathrm{PIP3}+\mathrm{PKBA}_c \to \mathrm{PKBA}_m+\mathrm{PIP3}$ & $k_{K,\mathrm{on}}[\mathrm{PIP3}][\mathrm{PKBA}_c]$ \\
4.11 & $\mathrm{PKBA}_m \to \mathrm{PKBA}_c$ & $k_{K,\mathrm{off}}[\mathrm{PKBA}_m]$ \\
4.12 & $\mathrm{PKBA}_m+\mathrm{TORC2}^* \to \mathrm{PKBA}_m^*+\mathrm{TORC2}^*$ & $k_{K,\mathrm{act}}[\mathrm{PKBA}_m][\mathrm{TORC2}^*]$ \\
4.13 & $\mathrm{PKBA}_m^* \to \mathrm{PKBA}_c$ & $k_{K,\mathrm{deact}}[\mathrm{PKBA}_m^*]$ \\
4.14 & $\mathrm{PKBA}_m^*+\mathrm{PI3K}_m \to \mathrm{PKBA}_m^*+\mathrm{PI3K}_c$ & $k_{K,\mathrm{fb}}[\mathrm{PKBA}_m^*][\mathrm{PI3K}_m]$ \\
\bottomrule
\end{tabularx}
\end{table*}

\begin{table*}[!t]
\centering
\caption{Module~4 constants. All rates except the PkbA arm (4.10--4.14) carry the uniform time-scale factor $s_4=20$ (Sec.~\ref{sec:m4}).}
\label{tab:m4const}
\footnotesize
\renewcommand{\arraystretch}{1.15}
\begin{tabularx}{\textwidth}{@{}l r l c >{\raggedright\arraybackslash}X@{}}
\toprule
Symbol & Value & Unit & Basis & Meaning, justification, source\\
\midrule
$s_4$ & \num{20} &  & C & Uniform clock multiplier on all Module~4 rates. Leaves every steady state and loop gain unchanged; it makes PIP3 formation complete within seconds of a 40\,s stimulus. \\
$[\mathrm{PI3K}]_\mathrm{tot}$ & \num{200000} & $\mathrm{molec/voxel}$ & E & Pool large enough that catalyst depletion is negligible. \\
$[\mathrm{PIP2}]_\mathrm{tot}$ & \num{769231} & $\mathrm{molec/voxel}$ & E & Lipid pool ($10^7$ per cell), conserved: $\mathrm{PIP2}+\mathrm{PIP3}=$ const. \\
$[\mathrm{PTEN}]_\mathrm{tot}$ & \num{10000} & $\mathrm{molec/voxel}$ & E & PTEN pool, conserved between cytosol and membrane. \\
$[\mathrm{PKBA}]_\mathrm{tot}$ & \num{100000} & $\mathrm{molec/voxel}$ & C & Pool of the PH-domain kinase PkbA; enlarged so that its feedback can act at low activated fractions. \\
$k_{3K,\mathrm{on}}$ & \num{0.4} & $\mu\mathrm{M}^{-1}\mathrm{s}^{-1}$ & C & \mkt{M4-ras-pi3k}{RasG-GTP recruits PI3K to the membrane}~\cite{sasaki2004,fukushima2019}. \\
$k_{3K,\mathrm{off}}$ & \num{10} & $\mathrm{s^{-1}}$ & C & PI3K membrane release; \mkt{M4-pip3-trans}{PI3K activity is transient after a stimulus}~\cite{huang2003}. \\
$k_\mathrm{syn}$ & \num{0.0955} & $\mu\mathrm{M}^{-1}\mathrm{s}^{-1}$ & C & PIP3 synthesis per membrane PI3K; sets the PIP3 count (shot noise), not its contrast. \\
$k_{PT,\mathrm{on}},\ k_{PT,\mathrm{off}}$ & \num{0.6},\ \num{4} & $\mu\mathrm{M}^{-1}\mathrm{s}^{-1},\ \mathrm{s^{-1}}$ & C & PTEN docking on PIP2 and release; \mkt{M4-recip}{PTEN sits on the membrane at the back of the cell}~\cite{funamoto2002,iijima2002}. \\
$k_\mathrm{ev}$ & \num{800} & $\mathrm{s^{-1}}$ & C & Saturated PIP3-dependent eviction rate of PTEN. \\
$K_\mathrm{ev},\ n_\mathrm{ev}$ & \num{1.31e5},\ 2 & $\mathrm{molec/voxel},\ -$ & C & Half-eviction level ($17\,\%$ of the lipid pool) and steepness of the eviction gate (Sec.~\ref{sec:hill}). \\
$k_{PT}$ & \num{1.39} & $\mu\mathrm{M}^{-1}\mathrm{s}^{-1}$ & C & \mkt{M4-pten-dephos}{PTEN dephosphorylation of PIP3 to PIP2}~\cite{iijima2002,huang2003}. \\
$k_\mathrm{SH}$ & \num{1.28} & $\mathrm{s^{-1}}$ & C & PTEN-independent 5-phosphatase removal of PIP3; $k_\mathrm{SH}=k_\mathrm{ship}[\mathrm{SHIP}]$ with a constant catalyst pool of 77 molecules/voxel. \\
$k_{K,\mathrm{on}}$ & \num{1.32e-3} & $\mu\mathrm{M}^{-1}\mathrm{s}^{-1}$ & D & PH-domain recruitment of PkbA; solved so that the resting rate is $0.002\,\mathrm{s^{-1}}$ at $7298$ PIP3/voxel. \\
$k_{K,\mathrm{off}}$ & \num{0.05} & $\mathrm{s^{-1}}$ & E & PkbA membrane release. \\
$k_{K,\mathrm{act}}$ & \num{0.507} & $\mu\mathrm{M}^{-1}\mathrm{s}^{-1}$ & D & TORC2-dependent activation; solved so that the resting rate is $8.5\times10^{-3}\,\mathrm{s^{-1}}$ at 81 TORC2$^*$/voxel. \\
$k_{K,\mathrm{deact}}$ & \num{0.02} & $\mathrm{s^{-1}}$ & E & PkbA$^*$ decay ($\tau=50$\,s); sets the refractory duration. \\
$k_{K,\mathrm{fb}}$ & \num{160} & $\mu\mathrm{M}^{-1}\mathrm{s}^{-1}$ & C & Strength of the PKB feedback on PI3K, chosen so that PIP3 returns to baseline after a pulse. \\
$D_{\mathrm{PI3K}_c},D_{\mathrm{PTEN}_c},D_{\mathrm{PKBA}_c}$ & \num{10} & $\mu\mathrm{m^2/s}$ & E & Cytosolic value; the PTEN shuttle crosses the cell in $\approx\!2.5$\,s. \\
$D_{\mathrm{PI3K}_m},D_{\mathrm{PTEN}_m},D_{\mathrm{PIP2/3}},D_{\mathrm{PKBA}_m}$ & \num{0.1} & $\mu\mathrm{m^2/s}$ & E & Membrane-bound species; the active PkbA is kept local so that feedback ends the patch where it forms. \\
\bottomrule
\end{tabularx}
\end{table*}


\paragraph{Reactions.}
\begin{description}\itemsep0pt
\item[4.1/4.2] RasG-GTP recruits cytosolic PI3K to the membrane, and membrane PI3K is released again. The release rate sets \mk{M4-pip3-trans}{how fast PI3K activity declines after the stimulus}~\cite{huang2003}.
\item[4.3] Membrane PI3K converts PIP2 to PIP3. The flux is proportional to the membrane PI3K count and does not depend on PIP3, so the rate constant sets the PIP3 count and not its contrast.
\item[4.4/4.5] PTEN docks on membrane PIP2 and is released; this gives PTEN its steady membrane fraction.
\item[4.6c] PIP3 evicts PTEN from the membrane through a cooperative gate; this is the \mk{M4-mutual}{mutual inhibition PIP3 $\dashv$ PTEN}~\cite{matsuoka2018}. The Hill-type form is used, instead of a mass-action form, because it separates a low-PIP3 state (PTEN on the membrane) from a high-PIP3 state (PTEN off the membrane).
\item[4.7] Membrane PTEN \mk{M4-pten-dephos}{removes the 3-phosphate of PIP3}~\cite{iijima2002,huang2003}.
\item[4.8] A PTEN-independent phosphatase (a 5-phosphatase such as SHIP) removes PIP3. It is the dominant sink at rest: without it PTEN, which is evicted by PIP3, would be the only route out of PIP3 and the resting cell would convert 60\,\% of the lipid pool to PIP3. The two-step dephosphorylation via PI(3,4)P$_2$ is collapsed to one step, which preserves the flux out of PIP3 but does not resolve PI(3,4)P$_2$, a species that is itself \mk{M4-pi34p2}{part of the Ras excitable network}~\cite{li2018}.
\item[4.10/4.11] PIP3 recruits the kinase PkbA (Akt) by its PH domain, and PkbA is released from the membrane. \mk{M4-pkba-recr}{PkbA is recruited to the membrane in response to cAMP in a PI3K-dependent way and its activation is transient}~\cite{meili1999}.
\item[4.12/4.13] Membrane PkbA is activated by TORC2, \mk{M4-tor-hm}{which is active at the front and acts through the hydrophobic motif}~\cite{kamimura2008}, and inactivated with $\tau=50$\,s. The two-step activation (PIP3-dependent recruitment first, TORC2-dependent phosphorylation second) is a coincidence detector and provides the lag behind PIP3 that a delayed negative feedback needs. The phosphoinositide-dependent kinases PdkA and PdkB, \mk{M4-pdk}{which phosphorylate the activation loop}~\cite{kamimura2010}, are not modelled; activation-loop and hydrophobic-motif phosphorylation are lumped into one step.
\item[4.14] Active PkbA removes PI3K from the membrane. \mk{M4-sca1fb}{The literature target of the feedback is the Sca1/RasGEF/PP2A complex and thus RasC}~\cite{charest2010}; because that complex is not resolved, the feedback is applied at its Module~4 end point (less Ras-GTP implies less membrane PI3K). Sign, locality and delay follow the published mechanism, and only the intermediate is omitted.
\end{description}
\mk{M4-two-pkb}{PkbA (with a PH domain, recruited by PIP3) and PKBR1 (myristoylated, constitutively at the membrane, PI3K-independent) are two different gene products}~\cite{meili1999,meili2000}; Module~4 owns PkbA and Module~8 owns PKBR1, so that no reaction describes the same species twice.

\paragraph{Constants.}
All rates of reactions 4.1--4.8 are multiplied by the clock factor $s_4=20$.
A uniform factor on all rates in a module is a clock, not a re-parameterisation: every steady state, every activated fraction and the loop gain are unchanged, and only the time scales are divided.
What limits PIP3 formation is the flux against the pool that has to be converted (about $1.7\times10^3$ molecules per voxel per second against a pool of about $6\times10^4$), and the factor was chosen so that PIP3 forms within seconds of a 40\,s stimulus.
The constants $k_\mathrm{syn}$, $k_{PT,\mathrm{on}}$ and $k_{PT}$ are expressed relative to legacy pool sizes (divided by the ratio of the present pool to the legacy pool), so that the flux per pool does not change when the pools were enlarged.
Enlarging a production-rate constant raises resting and peak PIP3 together, so it changes the PIP3 count (and hence the noise), not the contrast; it was therefore used as a noise lever and not as a contrast lever.
$k_{K,\mathrm{on}}$ and $k_{K,\mathrm{act}}$ are solved from resting-rate anchors (0.002 and $8.5\times10^{-3}$\,s$^{-1}$), so that the resting rate of the feedback arm is fixed and does not silently rescale when Modules~3 and~8 are retuned.
The slow recruitment ($k_{K,\mathrm{on}}$) delays the feedback, so that it acts after the cAMP pulse and not during it.
The PIP3 gate $K_\mathrm{ev}$ was set on the steep flank of the eviction curve; a value of the half-point above 19.5\,\% of the lipid pool made the PIP3 domain a winner-take-all latch whose orientation was set by the initial condition and not by the gradient.

\paragraph{Limitations.}
The mutual inhibition of Module~4 alone does not amplify a spatial signal: for the pools used its loop gain stays below one, so that the amplification of the polarity signal is provided by the Ras$\to$PIP3 feedback 3.3b and not by PTEN.
The PkbA feedback is required for a reset after a pulse: without it PIP3 keeps rising after the stimulus has ended.
The gate of reaction 4.12 reads TORC2, which is \mk{M4-tor-ras}{activated by RasC in the literature}~\cite{cai2010}; the coupling of Module~8 to Module~4 through TORC2$^*$ follows this.

\subsection{Module 8: RasC--TORC2--PKB--CRAC Relay}\label{sec:m8}

\paragraph{Role.}
The relay of the cAMP signal, i.e.\ the re-synthesis of cAMP in response to cAMP, runs through a second Ras branch.
RasC and RasG are both activated by cAMP, but RasC is responsible for adenylyl cyclase activation and RasG for chemotaxis; RasGEFA activates RasC but not RasG~\cite{kae2007}.
RasC is an upstream regulator of the target of rapamycin complex 2 (TORC2), which phosphorylates the myristoylated kinase PKBR1 independently of PIP3 and prolongs the time course of ACA activation~\cite{cai2010}.
Within seconds of a stimulus, TORC2 activates two PKB homologues, PkbA and PKBR1, also in cells lacking PI3K activity~\cite{kamimura2008}.
PKBR1 has no PH domain, is myristoylated for constitutive membrane localisation and does not require PI3K~\cite{meili2000}.
PKB and PKBR1 phosphorylate the scaffold Sca1 of a Sca1/RasGEF/PP2A complex and thereby control RasC in a negative feedback~\cite{charest2010}, and PKA exerts a further negative feedback on Ras, Rap1 and TORC2 signalling~\cite{scavello2017}.
The effector CRAC (a PH-domain protein) is required for receptor- and G-protein-mediated activation of ACA~\cite{insall1994}; CRAC binds PI3K products~\cite{comer2005} and is the translocation reporter of the G-protein pathway~\cite{parent1998}.
The model connects this chain to the excitable element of Module~0 by an output edge (8.14) that \emph{injects the excitable activator X}, and not by a direct activation of ACA.
A direct receptor-to-ACA edge would make the receptor arm a second pulse generator with its own amplitude and refractory state, whereas injecting X makes a relayed pulse identical to a spontaneous one and lets one refractory clock (Module~0) gate both.

\begin{table*}[!t]
\centering
\caption{Module~8 reactions: the RasC$\to$TORC2$\to$PKB$\to$CRAC relay arm and its output edge 8.14 into the excitable element of Module~0. PKBR1 (module 8b) and CRAC are the kinase and effector; $g(\mathrm{PIP3})=[1+(K_P/N_\mathrm{PIP3})^{8}]^{-1}$ and $b(\mathrm{Rf})$ are the gates of Sec.~\ref{sec:hill}.}
\label{tab:m8rxn}
\footnotesize
\renewcommand{\arraystretch}{1.15}
\begin{tabularx}{\textwidth}{@{}l>{\raggedright\arraybackslash}p{0.37\textwidth}>{\raggedright\arraybackslash}X@{}}
\toprule
ID & Reaction & Propensity (rate law)\\
\midrule
8.1 & $\mathrm{G}_{\beta\gamma}+\mathrm{GEF_A} \to \mathrm{GEF_A^*}+\mathrm{G}_{\beta\gamma}$ & $k_{A,\mathrm{on}}[\mathrm{G}_{\beta\gamma}][\mathrm{GEF_A}]$ \\
8.2 & $\mathrm{GEF_A^*} \to \mathrm{GEF_A}$ & $k_{A,\mathrm{off}}[\mathrm{GEF_A^*}]$ \\
8.2s & $\mathrm{B}^*+\mathrm{GEF_A^*} \to \mathrm{GEF_A}+\mathrm{B}^*$ & $k_S[\mathrm{B}^*][\mathrm{GEF_A^*}]$ \\
8.3 & $\mathrm{GEF_A^*}+\mathrm{RasC}^{\mathrm{GDP}} \to \mathrm{RasC}^{\mathrm{GTP}}+\mathrm{GEF_A^*}$ & $k_{C,\mathrm{on}}[\mathrm{GEF_A^*}][\mathrm{RasC}^{\mathrm{GDP}}]$ \\
8.4 & $\mathrm{RasC}^{\mathrm{GTP}} \to \mathrm{RasC}^{\mathrm{GDP}}$ & $k_{C,\mathrm{off}}[\mathrm{RasC}^{\mathrm{GTP}}]$ \\
8.4b & $\mathrm{RC}+\mathrm{CGAP} \to \mathrm{RC}+\mathrm{CGAP^*}$ & $k_{CG,\mathrm{on}}[\mathrm{RC}][\mathrm{CGAP}]$ \\
8.4c & $\mathrm{CGAP^*} \to \mathrm{CGAP}$ & $k_{CG,\mathrm{off}}[\mathrm{CGAP^*}]$ \\
8.4d & $\mathrm{CGAP^*}+\mathrm{RasC}^{\mathrm{GTP}} \to \mathrm{RasC}^{\mathrm{GDP}}+\mathrm{CGAP^*}$ & $k_{CG,\mathrm{hyd}}[\mathrm{CGAP^*}][\mathrm{RasC}^{\mathrm{GTP}}]$ \\
8.5 & $\mathrm{RasC}^{\mathrm{GTP}}+\mathrm{TORC2} \to \mathrm{TORC2^*}+\mathrm{RasC}^{\mathrm{GTP}}$ & $k_{T,\mathrm{on}}[\mathrm{RasC}^{\mathrm{GTP}}][\mathrm{TORC2}]$ \\
8.6 & $\mathrm{TORC2^*} \to \mathrm{TORC2}$ & $k_{T,\mathrm{off}}[\mathrm{TORC2^*}]$ \\
8b.1 & $\mathrm{TORC2^*}+\mathrm{PKBR1} \to \mathrm{PKBR1^*}+\mathrm{TORC2^*}$ & $k_{P,\mathrm{on}}[\mathrm{TORC2^*}][\mathrm{PKBR1}]$ \\
8b.2 & $\mathrm{PKBR1^*} \to \mathrm{PKBR1}$ & $k_{P,\mathrm{off}}[\mathrm{PKBR1^*}]$ \\
8b.3 & $\mathrm{PKBR1^*}+\mathrm{CRAC}_m \to \mathrm{CRAC}_m^*+\mathrm{PKBR1^*}$ & $k_\mathrm{Cr}[\mathrm{PKBR1^*}][\mathrm{CRAC}_m]$ \\
8b.4 & $\mathrm{PKA^*}+\mathrm{PKBR1^*} \to \mathrm{PKA^*}+\mathrm{PKBR1}$ & $k_{PKA}[\mathrm{PKA^*}][\mathrm{PKBR1^*}]$ \\
8.9 & $\mathrm{PKBA}_m^*+\mathrm{CRAC}_m \to \mathrm{CRAC}_m^*+\mathrm{PKBA}_m^*$ & $k_\mathrm{Cr}[\mathrm{PKBA}_m^*][\mathrm{CRAC}_m]$ \\
8.9b & $\mathrm{PIP3}+\mathrm{CRAC}_c \to \mathrm{CRAC}_m+\mathrm{PIP3}$ & $k_{\mathrm{Cr},P}[\mathrm{PIP3}][\mathrm{CRAC}_c]$ \\
8.9c & $\mathrm{CRAC}_c \to \mathrm{CRAC}_m$ & $k_{\mathrm{Cr},b}[\mathrm{CRAC}_c]$ \\
8.10 & $\mathrm{CRAC}_m^* \to \mathrm{CRAC}_m$ & $k_{\mathrm{Cr},\mathrm{dp}}[\mathrm{CRAC}_m^*]$ \\
8.10b & $\mathrm{CRAC}_m \to \mathrm{CRAC}_c$ & $k_{\mathrm{Cr},\mathrm{off}}[\mathrm{CRAC}_m]$ \\
8.14 & $\mathrm{CRAC}_m^* \to \mathrm{CRAC}_m^*+\mathrm{X}$ & $\dfrac{k_{X}}{N_\mathrm{CRAC}}\,N_{\mathrm{CRAC}^*}\,g(\mathrm{PIP3})\,b(\mathrm{Rf})\quad[\text{molec\,s}^{-1}\text{cell}^{-1}]$ \\
\bottomrule
\end{tabularx}
\end{table*}

\begin{table*}[!t]
\centering
\caption{Module~8 constants (a): RasC branch, TORC2 and PKBR1.}
\label{tab:m8const}
\footnotesize
\renewcommand{\arraystretch}{1.15}
\begin{tabularx}{\textwidth}{@{}l r l c >{\raggedright\arraybackslash}X@{}}
\toprule
Symbol & Value & Unit & Basis & Meaning, justification, source\\
\midrule
$[\mathrm{GEF_A}]_\mathrm{tot}$ & \num{1600} & $\mathrm{molec/voxel}$ & D & Sized so that the number of active RasGEFA at the driven state (67\,000 G$\beta\gamma$/voxel) equals that of the previous, saturated parameterisation. \\
$k_{A,\mathrm{on}}$ & \num{7.6e-4} & $\mu\mathrm{M}^{-1}\mathrm{s}^{-1}$ & D & $k_{E,\mathrm{on}}k_{A,\mathrm{off}}/k_{E,\mathrm{off}}$: RasGEFA and RasGEFR are both G$\beta\gamma$-driven GEFs~\cite{kae2007} and share the activated-fraction curve in G$\beta\gamma$. \\
$k_{A,\mathrm{off}}$ & \num{0.1} & $\mathrm{s^{-1}}$ & E & RasGEFA inactivation. \\
$k_S$ & \num{1} & $\mu\mathrm{M}^{-1}\mathrm{s}^{-1}$ & C & Release of active RasGEFA by the PKB-phosphorylated Sca1-like brake (negative feedback on RasC); calibrated so that RasC returns near baseline within $\approx\!40$\,s~\cite{charest2010}. \\
$k_{C,\mathrm{on}}$ & \num{2} & $\mu\mathrm{M}^{-1}\mathrm{s}^{-1}$ & E & RasGEFA-catalysed exchange on RasC~\cite{kae2007}. \\
$k_{C,\mathrm{off}}$ & \num{0.02} & $\mathrm{s^{-1}}$ & E & Intrinsic RasC hydrolysis (constitutive, slow). \\
$k_{CG,\mathrm{on}},\ k_{CG,\mathrm{off}},\ k_{CG,\mathrm{hyd}}$ & \num{0.05},\ \num{0.005},\ \num{150} & $\mu\mathrm{M}^{-1}\mathrm{s}^{-1},\ \mathrm{s^{-1}},\ \mu\mathrm{M}^{-1}\mathrm{s}^{-1}$ & E & Receptor-driven, slowly decaying RasC GAP that holds RasC adapted. \\
$k_{T,\mathrm{on}}$ & \num{50} & $\mu\mathrm{M}^{-1}\mathrm{s}^{-1}$ & C & TORC2 activation by RasC-GTP~\cite{cai2010}. \\
$k_{T,\mathrm{off}}$ & \num{0.5} & $\mathrm{s^{-1}}$ & C & TORC2 inactivation ($\tau=2$\,s). $k_{T,\mathrm{on}}$ and $k_{T,\mathrm{off}}$ were scaled together, which leaves the activated fraction unchanged and shortens only the latency, to a peak at 5--10\,s as observed for phosphorylation of Sca1~\cite{charest2010}. Other PKB substrates and the PKBR1 hydrophobic motif peak later, at 20--60\,s~\cite{kamimura2008,cai2010}, so the model is on the fast side. \\
$k_{P,\mathrm{on}},\ k_{P,\mathrm{off}}$ & \num{30},\ \num{0.5} & $\mu\mathrm{M}^{-1}\mathrm{s}^{-1},\ \mathrm{s^{-1}}$ & E & PKBR1 activation by TORC2 and its dephosphorylation ($\tau=2$\,s); PKBR1 is constitutively membrane-bound and PI3K-independent~\cite{meili2000,kamimura2008}. \emph{Departure:} the measured half-life of PKBR1 hydrophobic-motif phosphorylation is about 40\,s~\cite{cai2010}. \\
\bottomrule
\end{tabularx}
\end{table*}

\begin{table*}[!t]
\centering
\caption{Module~8 constants (b): CRAC, the relay output edge, pools and diffusion constants.}
\label{tab:m8constb}
\footnotesize
\renewcommand{\arraystretch}{1.15}
\begin{tabularx}{\textwidth}{@{}l r l c >{\raggedright\arraybackslash}X@{}}
\toprule
Symbol & Value & Unit & Basis & Meaning, justification, source\\
\midrule
$k_\mathrm{Cr}$ & \num{10} & $\mu\mathrm{M}^{-1}\mathrm{s}^{-1}$ & C & Phosphorylation of membrane-bound CRAC by either PKB. \mkt{M8-tor-pkbs}{Both PKBs are TORC2 targets}~\cite{kamimura2008} and \mkt{M8-aca-prolong}{RasC--TORC2--PKBR1 signalling prolongs ACA activation}~\cite{cai2010}; PKB$\to$CRAC is a model assumption. \\
$k_{\mathrm{Cr},P}$ & \num{1} & $\mu\mathrm{M}^{-1}\mathrm{s}^{-1}$ & E & PIP3-facilitated recruitment of CRAC, which is \mkt{M8-crac-pip3}{a PH-domain protein}~\cite{insall1994,parent1998,comer2005}. \\
$k_{\mathrm{Cr},b}$ & \num{0.2} & $\mathrm{s^{-1}}$ & C & PIP3-independent membrane association; keeps the relay functional when PIP3 is absent (\mkt{M8-pi3k-ind}{PI3K-independent PKB signalling}~\cite{kamimura2008}). \\
$k_{\mathrm{Cr},\mathrm{off}}$ & \num{0.1} & $\mathrm{s^{-1}}$ & E & Release of unphosphorylated CRAC from the membrane. \\
$k_{\mathrm{Cr},\mathrm{dp}}$ & \num{1.5} & $\mathrm{s^{-1}}$ & D & $=k_\mathrm{Cr}\,a\,(1-f)/f$ with $a=567$ PKBA$^*$/voxel (the measured peak) and target CRAC$^*$ fraction $f=0.44$; fixes the activated fraction of the CRAC switch. \\
$k_X$ & \num{200} & $\mathrm{molec\,s^{-1}\,cell^{-1}}$ & C & Injection rate of X at full CRAC$^*$ (per-CRAC$^*$ rate $k_X/N_\mathrm{CRAC}=200/2002=0.1\,\mathrm{s^{-1}}$, with $N_\mathrm{CRAC}=2002$ per cell); must exceed the $\approx\!24$-molecule threshold of Module~0. \\
$K_P,\ n_P$ & \num{5000},\ 8 & $\mathrm{molec/voxel},\ -$ & C & Half-open level and steepness of the PIP3 gate (Sec.~\ref{sec:hill}). \\
$k_\mathrm{PKA}$ & \num{0.938} & $\mu\mathrm{M}^{-1}\mathrm{s}^{-1}$ & E & PKA$^*$ feedback on PKBR1$^*$; \mkt{M8-pka-fb}{PKA exerts negative feedback on Ras, Rap1 and TORC2 signalling}~\cite{scavello2017}. The magnitude is a design choice: the value $0.6\,\mathrm{min^{-1}}$ is \mkt{M8-laub-k}{the coefficient $k_5$ of}~\cite{laub1998} (CAR1$\to$ERK2, not a PKA coefficient; the PKA$\to$ERK2 coefficient is $k_6=0.8\,\mathrm{min^{-1}}$, used in Module~10), $30\times$ boosted and normalised by the PKA pool. \\
$[\mathrm{RasC}]_\mathrm{tot}$ & \num{385} & $\mathrm{molec/voxel}$ & E & Pool size; not constrained by measurements. \\
$[\mathrm{CGAP}]_\mathrm{tot},\ [\mathrm{TORC2}]_\mathrm{tot},\ [\mathrm{CRAC}]_\mathrm{tot}$ & \num{154} & $\mathrm{molec/voxel}$ & E & Pool sizes (2\,002 per cell); not constrained by measurements. \\
$[\mathrm{PKBR1}]_\mathrm{tot}$ & \num{77} & $\mathrm{molec/voxel}$ & E & PKBR1 pool; PkbA (Module~4) is a separate, larger pool. \\
$D_{\mathrm{RasC}},D_{\mathrm{TORC2}},D_{\mathrm{PKBR1}},D_{\mathrm{CRAC}_m}$ & \num{0.1} & $\mu\mathrm{m^2/s}$ & E & Membrane-bound species. \\
$D_{\mathrm{GEF_A}},D_{\mathrm{CGAP}},D_{\mathrm{CRAC}_c}$ & \num{10} & $\mu\mathrm{m^2/s}$ & E & Cytosolic value ($D_{\mathrm{GEF_A^*}}=2$). \\
\bottomrule
\end{tabularx}
\end{table*}


\paragraph{Reactions.}
\begin{description}\itemsep0pt
\item[8.1/8.2] RasGEFA is activated by G$\beta\gamma$ and inactivated with $\tau=10$\,s. It uses the same G$\beta\gamma$-dependent activated-fraction curve as the RasG excitor 3.1, since both are G$\beta\gamma$-driven RasGEFs~\cite{kae2007}.
\item[8.2s] The PKB-phosphorylated Sca1-like brake substrate B$^*$ (Module~3) releases active RasGEFA, which is the negative feedback of~\cite{charest2010}. Without PKB or TORC2, RasC activation is increased and fails to adapt by 40\,s; the constant $k_S$ is calibrated to return RasC near baseline within about 40\,s.
\item[8.3/8.4] Active RasGEFA converts RasC-GDP to RasC-GTP; RasC-GTP hydrolyses slowly.
\item[8.4b--d] A receptor-driven, slowly decaying RasC GAP (RC $\to$ CGAP$^*$, decay $\tau=200$\,s) hydrolyses RasC-GTP and holds RasC in its adapted state.
\item[8.5/8.6] RasC-GTP activates TORC2~\cite{cai2010}, which relaxes with $\tau=2$\,s. The short lifetime is a design choice that places the peak of TORC2 targets within 5--10\,s, as observed for the PKB substrate Sca1~\cite{charest2010}.
\item[8b.1/8b.2] TORC2$^*$ phosphorylates the membrane-anchored PKBR1 directly (no recruitment step), and PKBR1$^*$ is dephosphorylated rapidly, so that PKBR1 is a fast-on/fast-off pulse generator that is cleared before the slow negative feedbacks accumulate. Measured PKBR1 hydrophobic-motif phosphorylation is slower: it peaks at 30--60\,s, has a half-life of about 40\,s and is back at baseline after 2--3\,min~\cite{cai2010}; the model's $\tau=2$\,s is a design choice for the fast relay, not a measured value.
\item[8b.4] PKA$^*$ inactivates PKBR1$^*$, which closes a slow negative feedback from cAMP to the relay arm~\cite{scavello2017}.
\item[8.9, 8b.3] Both PKBs phosphorylate the membrane-bound CRAC pool, PkbA (Module~4) through 8.9 and PKBR1 through 8b.3.
\item[8.9b/8.9c/8.10/8.10b] CRAC is a three-state moiety: cytosolic CRAC$_c$, membrane-bound CRAC$_m$ and phosphorylated CRAC$_m^*$. Two facts must hold together and are compatible only if recruitment and activation are separate steps: CRAC is a PH-domain protein recruited by PIP3~\cite{parent1998,comer2005}, and PKB signalling and the relay do not require PIP3~\cite{kamimura2008}. PIP3 therefore recruits CRAC (8.9b), a PIP3-independent association keeps the relay functional without PIP3 (8.9c), and the kinases phosphorylate CRAC on the membrane (8.9, 8b.3). Dephosphorylation (8.10) returns CRAC$^*$ to the membrane pool and not to the cytosol, because a phosphatase does not strip the lipid anchor; releasing it to the cytosol would cap the activated fraction at $k_\mathrm{rec}/(k_\mathrm{rec}+k_\mathrm{off})$, whatever the kinase drive. Release from the membrane (8.10b) is a separate step.
\item[8.14] CRAC$^*$ injects X into the excitable element at up to $k_X=200$ molecules\,s$^{-1}$ per cell, gated by PIP3 and by the refractory clock. The PIP3 gate requires that the cell has already resolved the direction of the signal before it re-amplifies it: PIP3 is the first species of the cascade that is both stimulus-driven and polarised.
\end{description}
Three further edges exist in the code and are excluded from the default network:
the activation of PkbA by TORC2 without a recruitment step (8.7/8.8) is PKBR1's mechanism applied to PkbA's species and would double count the kinase that Module~8b already provides;
the PKA feedback on PkbA (8.13) is excluded together with them;
and the direct activation of ACA by CRAC$^*$ (8.11) is disabled ($k=0$) for the reason stated above.

\paragraph{Constants.}
Most pool sizes and the RasC-branch rate constants are estimates for which no measurements exist; the calibrated properties are the timing (TORC2$^*$ and PKBR1$^*$ peak within 5--10\,s), the activated fraction of the CRAC switch and the relay drive.
Two constants are derived from measurements of other modules and must be re-derived when those change.
First, $k_{\mathrm{Cr},\mathrm{dp}}$ sets the activated fraction of the CRAC phospho-cycle, which is a function of $a/k_{\mathrm{Cr},\mathrm{dp}}$ only, where $a=k_\mathrm{Cr}[\mathrm{PKBA}^*]$; it is derived from the measured peak of PkbA$^*$ (567 molecules per voxel) and the target CRAC$^*$ fraction of 44\,\%.
When Module~4 was re-anchored the measured peak fell tenfold, the fraction fell to 11\,\% and the relay stopped firing, so the derivation is stored with the measurement it depends on.
Second, the half-open level $K_P$ of the PIP3 gate is placed between the PIP3 peak evoked by 2\,nM (about 5000 molecules per voxel) and the standing PIP3 (about $10^3$).
Neither the gain $k_X$ nor $k_\mathrm{Cr}$ can replace these derivations: both would restore the product of drive and fraction but leave a switch operating in a tenth of its range, without its contrast.
The pairing $k_{A,\mathrm{on}}=k_{E,\mathrm{on}}k_{A,\mathrm{off}}/k_{E,\mathrm{off}}$ removes a saturation that the earlier parameterisation had (RasGEFA 83\,\% active at rest and 99\,\% at 60\,nM), where a cAMP step drove RasC-GTP down instead of up; RasC is activated rapidly and transiently by cAMP~\cite{kae2004}.
The pool $[\mathrm{GEF_A}]_\mathrm{tot}$ was chosen so that the number of active GEFs at the driven state is preserved.
The relay is nearly dose independent in the model: 2\,nM cAMP fires the excitable element as fully as 200\,nM, because the LEGI/PIP3 stage reports contrast rather than level; this is why the gate must be steep.

\subsection{Module 9: cAMP Synthesis and Export}\label{sec:m9}

\paragraph{Role.}
Cells lacking the aggregation-stage adenylyl cyclase ACA have almost no cyclase activity and fail to aggregate, whereas motility and chemotaxis are unaffected, so cAMP synthesis is what allows cell--cell communication~\cite{pitt1992}.
ACA is activated in response to cAMP through CRAC~\cite{insall1994} and the MAP kinase ERK2~\cite{segall1995}; after a sustained stimulus its activity peaks after about 30\,s and returns to baseline with a half-time of 1.8\,min~\cite{cai2010}.
cAMP is exported by ABC transporters, of which AbcB3 is a component~\cite{miranda2015}, and the rate of secretion is proportional to the intracellular cAMP concentration~\cite{dinauer1980}.
Cells released about 5\,pmol cAMP per $10^6$ cells, i.e.\ about $3\times10^{6}$ molecules per cell, in response to a 2-min stimulus at 7\,h of development~\cite{devreotes1979}.

\begin{table*}[!t]
\centering
\caption{Module~9 reactions: cAMP synthesis by the adenylyl cyclase ACA, export, and ACA inactivation. Reaction 8.12 belongs to the ACA arm and is listed here. A coloured reaction is the one to which the equally coloured statement in the text and the equally coloured passage in the cited paper refer; a small square appends a further statement.}
\label{tab:m9rxn}
\footnotesize
\renewcommand{\arraystretch}{1.15}
\begin{tabularx}{\textwidth}{@{}l>{\raggedright\arraybackslash}p{0.37\textwidth}>{\raggedright\arraybackslash}X@{}}
\toprule
ID & Reaction & Propensity (rate law)\\
\midrule
\mkt{M9-aca-synth}{9.1} & \mkt{M9-aca-synth}{$\mathrm{ACA}^* \to \mathrm{ACA}^*+\mathrm{cAMP}_i$} & $k_{9}\,(1+\alpha_H H/H_{\max})[\mathrm{ACA}^*]$ \\
\mkt{M9-aca-synth}{9.2} & \mkt{M9-aca-synth}{$\mathrm{ACA}^* \to \mathrm{ACA}^*+\mathrm{cAMP}_i$} & $k_{11}[\mathrm{ACA}^*]$ \\
\mkn{M9-secr-prop}{M9-abcb3}{9.2b} & \mkn{M9-secr-prop}{M9-abcb3}{$\mathrm{cAMP}_i \to \mathrm{cAMP}_e$} & $k_{x}[\mathrm{cAMP}_i]$ \\
\mkn{M9-pka-null}{M9-laub-k2}{9.3} & \mkn{M9-pka-null}{M9-laub-k2}{$\mathrm{PKA}^*+\mathrm{ACA}^* \to \mathrm{PKA}^*+\mathrm{ACA}$} & $k_{2}'[\mathrm{PKA}^*][\mathrm{ACA}^*]$ \\
\mkt{M9-aca-kin}{8.12} & \mkt{M9-aca-kin}{$\mathrm{ACA}^* \to \mathrm{ACA}$} & $k_{Ao}[\mathrm{ACA}^*]$ \\
\bottomrule
\end{tabularx}
\end{table*}

\begin{table*}[!t]
\centering
\caption{Module~9 constants.}
\label{tab:m9const}
\footnotesize
\renewcommand{\arraystretch}{1.15}
\begin{tabularx}{\textwidth}{@{}l r l c >{\raggedright\arraybackslash}X@{}}
\toprule
Symbol & Value & Unit & Basis & Meaning, justification, source\\
\midrule
$k_9$ & \num{8} & $\mathrm{s^{-1}}$ & C & Hunger-scaled synthesis arm; $k_9+k_{11}=20\,\mathrm{s^{-1}}$ is an ordinary adenylyl-cyclase turnover chosen so that one pulse exports $\approx\!9\times10^6$ cAMP molecules per cell, three times the \mkt{M9-pulse-amp}{$3\times10^6$ measured for a 2\,min stimulus at 7\,h}~\cite{devreotes1979}. \\
$k_{11}$ & \num{12} & $\mathrm{s^{-1}}$ & C & Constant synthesis arm (see $k_9$). \\
$\alpha_H$ & \num{0.5} &  & E & Hunger amplification of synthesis ($1+\alpha_H H/H_{\max}$); starved cells make more cAMP per active ACA. \\
$k_x$ & \num{0.5} & $\mathrm{s^{-1}}$ & C & First-order export via \mkt{M9-abcb3}{the ABC transporter AbcB3}~\cite{miranda2015}: \mkt{M9-secr-prop}{the secretion rate is proportional to intracellular cAMP}~\cite{dinauer1980}. \emph{Departure:} measured first-order secretion constants are \mkt{M9-ks}{$0.34$ and $0.94\,\mathrm{min^{-1}}$ ($0.006$--$0.016\,\mathrm{s^{-1}}$)}~\cite{dinauer1980}, i.e.\ 30--90 times slower than the value used ($\tau=2$\,s), which was chosen so that export adds no delay to the $\approx\!40$\,s pulse and RegA (Module~10) controls the exported fraction. \\
$k_2'$ & \num{0.0469} & $\mu\mathrm{M}^{-1}\mathrm{s}^{-1}$ & E & PKA-dependent termination of ACA activation, a lumped edge: \mkt{M9-pka-null}{in the Laub--Loomis circuit PKA terminates ACA activation indirectly, by inhibiting ERK2 and the receptor, and PKA-null cells keep ACA active for much longer}~\cite{laub1998}. The magnitude equals \mkt{M9-laub-k2}{the coefficient $k_2=0.9\,\mathrm{min^{-1}}$ of}~\cite{laub1998}, which there is the \emph{spontaneous} ACA inactivation, divided by the PKA pool ($0.32\,\mu$M); it is a calibration, not a measurement. \\
$k_{Ao}$ & \num{6.4e-3} & $\mathrm{s^{-1}}$ & L & Intrinsic ACA$^*$ decay, $t_{1/2}=108$\,s, \mkt{M9-aca-kin}{the decline of ACA activity after a sustained stimulus (peak at 30\,s, $t_{1/2}=1.8$\,min}~\cite{cai2010}); the realised decline is faster because 9.3 adds to it. \\
$[\mathrm{ACA}]_\mathrm{tot}$ & \num{865} & $\mathrm{molec/voxel}$ & C & See Module~0. \\
$D_{\mathrm{cAMP}_i}$ & \num{30} & $\mu\mathrm{m^2/s}$ & E & Buffered effective cytosolic diffusion of cAMP. \\
$D_{\mathrm{ACA}}$ & \num{0.1} & $\mu\mathrm{m^2/s}$ & E & Membrane-bound cyclase. \\
\bottomrule
\end{tabularx}
\end{table*}


\paragraph{Reactions.}
\begin{description}\itemsep0pt
\item[9.1/9.2] The active cyclase ACA$^*$ makes cAMP inside the cell. The synthesis is split into a Hunger-scaled arm and a constant arm, so that starved cells, which are closer to the oscillatory regime, make more cAMP per active cyclase; the two arms sum to $20$\,s$^{-1}$ per ACA$^*$.
\item[9.2b] cAMP is exported to the extracellular medium. The export is first order in the intracellular concentration~\cite{dinauer1980} and is the only cAMP output of the cell: removing it removes the secreted signal although the intracellular cascade still fires.
\item[9.3] PKA$^*$ inactivates ACA$^*$. In the Laub--Loomis circuit PKA terminates ACA activation only indirectly, through inhibition of ERK2 and of the receptor CAR1; in mutants lacking PKA, ACA and ERK2 rise normally but stay high for much longer~\cite{laub1998}. The direct edge is a lumped simplification of this negative feedback.
\item[8.12] ACA$^*$ decays spontaneously. The rate reproduces the decline of ACA activity after a sustained stimulus~\cite{cai2010}. At the much faster value used earlier, ACA$^*$ was slaved to the X spike and the pulse carried only $8\times10^5$ molecules; at the present value the brief X spike \emph{charges} ACA$^*$, and the pulse length is set by the ACA$^*$/PKA/RegA loop and not by X.
\end{description}

\paragraph{Constants.}
The amplitude is anchored to the measured relay output and not to the normalised synthesis coefficients of the Laub--Loomis model~\cite{laub1998}: the pool of 865 ACA molecules per voxel and the turnover of $20\,\mathrm{s^{-1}}$ were chosen so that one pulse exports about $9\times10^{6}$ molecules per cell.
This is three times the $3\times10^6$ molecules measured for a 2-min stimulus~\cite{devreotes1979}, so the modelled pulse is at the upper end of measured amplitudes.
The cyclase pool has never been measured.
The export rate $k_x$ was lowered from an earlier derived value of $3\,\mathrm{s^{-1}}$: at that value export outran the phosphodiesterase RegA by a factor of 250, and knocking out the whole ERK2$\to$RegA edge changed the exported cAMP by a factor of 1.0.
At $0.5$\,s$^{-1}$ the exported fraction is 0.05 at rest (RegA armed) and 0.83 at the pulse peak (RegA disarmed), so the contrast between rest and pulse is carried by RegA and not by the transporter.
This is much faster than what was measured in perfused amoebae: the first-order secretion constant was $0.34$--$0.94\,\mathrm{min^{-1}}$ and that of intracellular phosphodiesterase destruction about $1.7\,\mathrm{min^{-1}}$, so that 16\,\% (for a $10^{-6}$\,M stimulus) to 47\,\% (for $10^{-8}$\,M) of the newly made cAMP was secreted~\cite{dinauer1980}.
Both the export and the RegA rate of the model are therefore one to two orders of magnitude faster than these estimates, and the modelled exported fraction at the pulse peak (0.83) is above the measured range; the model buys a fast, sharp pulse (about 40\,s) with fast turnover. The measured pulse is broader, and the measured rate constants would need a correspondingly slower relay.
\emph{Limitation:} only a first-order dependence of export on cAMP$_i$ is established~\cite{dinauer1980}; its numerical value is a calibration.

\subsection{Module 10: RegA, PKA and ERK2 Brake}\label{sec:m10}

\paragraph{Role.}
The intracellular phosphodiesterase RegA (DdPDE2) \mk{M10-rega-pde}{degrades intracellular cAMP}~\cite{bader2007}.
It carries a phospho-relay receiver domain: \mk{M10-relay}{phosphorylation of Asp212 by the histidine phospho-transfer protein RdeA}~\cite{thomason1998,thomason1999} \mk{M10-rega-20}{activates the enzyme at least 20-fold}~\cite{thomason1999}, and \mk{M10-rega-km}{its cAMP-hydrolysing domain has a Michaelis constant of about 5\,$\mu$M}~\cite{thomason1998}.
\mk{M10-rega-loss}{Loss of RegA activates PKA and disturbs differentiation}~\cite{shaulsky1998}.
cAMP oscillations in populations are generated by a circuit in which \mk{M10-circuit}{cAMP activates PKA, PKA inhibits ERK2, and ERK2 controls cAMP levels through RegA}~\cite{laub1998,maeda2004}; \mk{M10-erk-phase}{the ERK2 activation oscillates in phase with the cAMP peaks}~\cite{laub1998,maeda2004}.
PKA also \mk{M10-pka-erk}{shapes ERK2 activation and adaptation}~\cite{aubry1997}.
Module~10 is the resulting brake: at rest RegA holds intracellular cAMP near zero, an ERK2 spike disarms it so that cAMP can accumulate, and its recovery and the PKA arm terminate the pulse and set the refractory tail.

\begin{table*}[!t]
\centering
\caption{Module~10 reactions: the intracellular brake. RegA is a conserved three-state moiety: RegA$_u$ (unphosphorylated, weakly active), RegA$_p$ (phosphorylated at D212, active) and RegA$_{pi}$ (phosphorylated by ERK2, inactive). PKA and ERK2 are the kinases of the Laub--Loomis/Maeda circuit.}
\label{tab:m10rxn}
\footnotesize
\renewcommand{\arraystretch}{1.15}
\begin{tabularx}{\textwidth}{@{}l>{\raggedright\arraybackslash}p{0.37\textwidth}>{\raggedright\arraybackslash}X@{}}
\toprule
ID & Reaction & Propensity (rate law)\\
\midrule
10.B1 & $\mathrm{RegA}_u \to \mathrm{RegA}_p$ & $k_{B1}[\mathrm{RegA}_u]$ \\
10.B2 & $\mathrm{RegA}_p \to \mathrm{RegA}_u$ & $k_{B2}[\mathrm{RegA}_p]$ \\
10.B3 & $\mathrm{RegA}_p+\mathrm{ERK2}^* \to \mathrm{RegA}_{pi}+\mathrm{ERK2}^*$ & $k_{B3}[\mathrm{RegA}_p][\mathrm{ERK2}^*]$ \\
10.B4 & $\mathrm{RegA}_{pi} \to \mathrm{RegA}_p$ & $k_{B4}[\mathrm{RegA}_{pi}]$ \\
10.C1 & $\mathrm{cAMP}_i+\mathrm{RegA}_p \to \mathrm{RegA}_p$ & $k_\mathrm{cat}[\mathrm{cAMP}_i][\mathrm{RegA}_p]$ \\
10.C2 & $\mathrm{cAMP}_i+\mathrm{RegA}_u \to \mathrm{RegA}_u$ & $\phi\,k_\mathrm{cat}[\mathrm{cAMP}_i][\mathrm{RegA}_u]$ \\
10.D1 & $\mathrm{cAMP}_i+\mathrm{PKA} \to \mathrm{cAMP}_i+\mathrm{PKA}^*$ & $k_{D1}[\mathrm{cAMP}_i][\mathrm{PKA}]$ \\
10.D2 & $\mathrm{PKA}^* \to \mathrm{PKA}$ & $k_{D2}[\mathrm{PKA}^*]$ \\
10.D3 & $\mathrm{PKA}^*+\mathrm{ERK2}^* \to \mathrm{PKA}^*+\mathrm{ERK2}$ & $k_{D3}[\mathrm{PKA}^*][\mathrm{ERK2}^*]$ \\
\bottomrule
\end{tabularx}
\end{table*}

\begin{table*}[!t]
\centering
\caption{Module~10 constants.}
\label{tab:m10const}
\footnotesize
\renewcommand{\arraystretch}{1.15}
\begin{tabularx}{\textwidth}{@{}l r l c >{\raggedright\arraybackslash}X@{}}
\toprule
Symbol & Value & Unit & Basis & Meaning, justification, source\\
\midrule
$[\mathrm{RegA}]_\mathrm{tot}$ & \num{1540} & $\mathrm{molec/voxel}$ & E & Copy number never measured; $0.32\,\mu$M puts RegA on the scale of ERK2 and PKA and keeps the required catalytic efficiency sub-diffusion-limited. \\
$[\mathrm{PKA}]_\mathrm{tot}$ & \num{1540} & $\mathrm{molec/voxel}$ & E & $0.32\,\mu$M. \\
$k_{B1}$ & \num{0.05} & $\mathrm{s^{-1}}$ & E & \mkt{M10-relay}{Lumped phospho-relay (RdeA$\to$RegA D212) that activates RegA}~\cite{thomason1999,thomason1998}. \\
$k_{B2}$ & \num{5.6e-3} & $\mathrm{s^{-1}}$ & C & \mkt{M10-relay}{Reverse phospho-transfer}~\cite{thomason1999}; chosen for a resting active fraction of 0.9. \\
$k_{B3}$ & \num{8} & $\mu\mathrm{M}^{-1}\mathrm{s}^{-1}$ & C & \mkt{M10-erk-rega}{ERK2$^*$ phosphorylates and thereby inhibits RegA}~\cite{laub1998,maeda2004}. The Laub--Loomis coefficient (\mkt{M10-laub-k8}{$k_8=1.3$ in relative units}~\cite{laub1998}) is not a molecular rate; the value gives $\tau\approx0.5$\,s at the ERK2$^*$ peak so that the brake is released inside the X spike. \\
$k_{B4}$ & \num{0.0167} & $\mathrm{s^{-1}}$ & C & The reset reaction ($\tau=60$\,s) that ends each pulse; of the order of the RegA replenishment coefficient \mkt{M10-laub-k7}{$k_7=2\,\mathrm{min^{-1}}$ of the Laub--Loomis model}~\cite{laub1998} and set so that the cell has reset within about 6\,min. \\
$k_\mathrm{cat}$ & \num{30} & $\mu\mathrm{M}^{-1}\mathrm{s}^{-1}$ & D & $k_\mathrm{cat}/K_M=3\times10^{7}\,\mathrm{M^{-1}s^{-1}}$: $k_\mathrm{cat}\approx150\,\mathrm{s^{-1}}$ at \mkt{M10-rega-km}{the measured $K_M\approx5\,\mu$M}~\cite{thomason1998}. \\
$\phi$ & \num{0.05} &  & L & Activity of RegA$_u$ relative to RegA$_p$; \mkt{M10-rega-20}{phosphorylation activates RegA by at least 20-fold}~\cite{thomason1999}. \\
$k_{D1}$ & \num{0.25} & $\mu\mathrm{M}^{-1}\mathrm{s}^{-1}$ & E & cAMP binding to the PKA regulatory subunit; with $k_{D2}$ this gives $K_d=k_{D2}/k_{D1}=0.1\,\mu$M (model choice; \mkt{M10-circuit}{PKA activation is the second element of the circuit of}~\cite{laub1998}). \\
$k_{D2}$ & \num{0.025} & $\mathrm{s^{-1}}$ & D & PKA$^*$ decay ($\tau=40$\,s), from \mkt{M10-laub-k4}{the Laub--Loomis coefficient $k_4=1.5\,\mathrm{min^{-1}}$}~\cite{laub1998}; sets the minutes-long refractory tail. \\
$k_{D3}$ & \num{0.0417} & $\mu\mathrm{M}^{-1}\mathrm{s}^{-1}$ & D & PKA$^*$ inhibits ERK2, from \mkt{M10-laub-k6}{the coefficient $k_6=0.8$ of}~\cite{laub1998} (relative units, normalised here by the PKA pool); PKA also \mkt{M10-pka-erk}{shapes ERK2 activation and adaptation}~\cite{aubry1997}. \\
$D_{\mathrm{RegA}},D_{\mathrm{PKA}},D_{\mathrm{ERK2}}$ & \num{10} & $\mu\mathrm{m^2/s}$ & E & Cytosolic value. \\
\bottomrule
\end{tabularx}
\end{table*}


\paragraph{Reactions.}
\begin{description}\itemsep0pt
\item[10.B1/10.B2] The phospho-relay activates RegA and reverses. The pair sets the resting active fraction, 0.9. Only \mk{M10-relay}{the lumped transfer from RdeA to RegA} is modelled~\cite{thomason1999}.
\item[10.B3] ERK2$^*$ converts active RegA into an inactive, ERK2-modified form. The model implements \mk{M10-erk-rega}{the control of RegA by ERK2}~\cite{maeda2004} as inactivation, the sign required for an ERK2 spike to escape the brake.
\item[10.B4] The inactive form returns to the active form ($\tau=60$\,s), since \mk{M10-rega-phosphatase}{a protein phosphatase activates RegA}~\cite{laub1998}. This reset step is the one that terminates each pulse; without it RegA would be lost irreversibly.
\item[10.C1/10.C2] RegA hydrolyses cAMP$_i$. The rate is linear in cAMP because cAMP$_i$ stays far below the Michaelis constant; \mk{M10-rega-20}{the unphosphorylated form keeps $1/20$ of the activity of the phosphorylated form}~\cite{thomason1999}.
\item[10.D1/10.D2] \mk{M10-pka-r}{cAMP activates PKA by binding to its regulatory subunit}~\cite{laub1998}, and PKA$^*$ decays with $\tau=40$\,s, which sets the minutes-long tail of the refractory period.
\item[10.D3] \mk{M10-circuit}{PKA$^*$ inactivates ERK2$^*$}, which closes the loop~\cite{laub1998,aubry1997}.
\end{description}

\paragraph{Constants.}
The Laub--Loomis coefficients are rates of a normalised-activity model and not molecular rate constants, so they cannot be used literally.
The Laub--Loomis coefficients refer to activities in relative units, so they cannot be used as molecular rate constants.
In the calibration of this model, a rate for RegA inactivation taken over literally kept the brake closed (knock-down $1.3\times$); the present value gives a knock-down of RegA$_p$ of $64\times$.
The cyclase-RegA competition needs $k_\mathrm{cat}/K_M\approx3\times10^{7}\,\mathrm{M^{-1}s^{-1}}$ with the RegA pool of $0.32\,\mu$M, which is sub-diffusion-limited; an earlier pool of $0.032\,\mu$M would have needed a catalytic efficiency near the diffusion limit before RegA could compete with export.
The copy numbers of RegA, PKA and ERK2 have not been measured.
Only $k_{D2}$ and $k_{D3}$ are tied to coefficients of the Laub--Loomis model ($k_4=1.5$ and $k_6=0.8$ in relative units, the latter divided by the PKA pool); $k_{B4}$ is calibrated, and the resulting pulse (an export of about $10^7$ molecules, a RegA knock-down of $64\times$ and a full reset within about 6\,min) is what the module is calibrated to.

\subsection{Module 11: Extracellular Phosphodiesterase PdsA and Inhibitor PdiA}\label{sec:m11}

\paragraph{Role.}
\mk{M11-pde-sec}{The level of extracellular cAMP is controlled by a phosphodiesterase that is secreted by the cells}~\cite{sucgang1997}; \mk{M11-pde-bader}{extracellular cAMP is degraded predominantly by DdPDE1 (PdsA) and its close homologue}~\cite{bader2007}.
\mk{M11-pde-null-g}{Cells lacking PdsA respond chemotactically to a narrower range of cAMP concentrations and do not form streams at low density}~\cite{garcia2009}, and \mk{M11-pde-null-s}{PdsA-null cells cannot move in a coordinated manner in mounds}~\cite{sucgang1997}.
The \emph{pdsA} gene has \mk{M11-promoters}{separate promoters for growth, aggregation and late development}~\cite{faure1990}, and \mk{M11-pa-off}{the aggregative promoter is switched off after aggregation}~\cite{weening2003}.
PdsA is inhibited by a secreted glycoprotein, PdiA, \mk{M11-inhib-stoich}{with 1:1 stoichiometry}~\cite{franke1981}; \mk{M11-inhib-absent}{the inhibitor is absent in growing cells and appears after transfer to starvation buffer}~\cite{franke1981}, and \mk{M11-inhib-reg}{its level is regulated independently by exogenous cAMP}~\cite{yeh1978}.
Module~11 therefore contains the synthesis and secretion chain of the enzyme, a starvation-induced inhibitor that is repressed by the cAMP history, and their reversible binding.

\begin{table*}[!t]
\centering
\caption{Module~11 reactions: the extracellular phosphodiesterase PdsA and its inhibitor PdiA. PdsA$_i$, PdsA$_m$ and PdsA$_e$ are the intracellular precursor, the membrane-displayed and the secreted enzyme; cAMP$_s$ is a slow, low-pass filtered copy of the local cAMP$_e$. Reactions marked $\ast$ run twice: inside cells and, with identical constants, on agent-free voxels of the world. $P(H)=H^2/(H_P^2+H^2)$ and $r(\mathrm{cAMP}_s)=[1+(\mathrm{cAMP}_s/K_\mathrm{pdi})^2]^{-1}$ are the Hill gates of Sec.~\ref{sec:hill}.}
\label{tab:m11rxn}
\footnotesize
\renewcommand{\arraystretch}{1.15}
\begin{tabularx}{\textwidth}{@{}l>{\raggedright\arraybackslash}p{0.37\textwidth}>{\raggedright\arraybackslash}X@{}}
\toprule
ID & Reaction & Propensity (rate law)\\
\midrule
11.D1a & $\varnothing \to \mathrm{PdsA}_i$ & $k_\mathrm{bas}\ \text{(zeroth order)}$ \\
11.D1b & $\varnothing \to \mathrm{PdsA}_i$ & $k_\mathrm{ind}\,P(H)\ \text{(zeroth order)}$ \\
11.D2 & $\mathrm{PdsA}_i \to \mathrm{PdsA}_m$ & $k_\mathrm{sec}[\mathrm{PdsA}_i]$ \\
11.D3 & $\mathrm{PdsA}_m \to \mathrm{PdsA}_e$ & $k_\mathrm{shed}[\mathrm{PdsA}_m]$ \\
11.D4a,b & $\mathrm{PdsA}_m \to \varnothing;\ \ \mathrm{PdsA}_i \to \varnothing$ & $k_\mathrm{deg}[\,\cdot\,]$ \\
11.D4c$^\ast$ & $\mathrm{PdsA}_e \to \varnothing$ & $k_\mathrm{deg}[\mathrm{PdsA}_e]$ \\
11.D4d$^\ast$ & $\mathrm{PdsA{:}PdiA} \to \varnothing$ & $k_\mathrm{deg}[\mathrm{PdsA{:}PdiA}]$ \\
11.F1 & $\mathrm{cAMP}_e \to \mathrm{cAMP}_e+\mathrm{cAMP}_s$ & $k_\mathrm{trk}[\mathrm{cAMP}_e]$ \\
11.F2 & $\mathrm{cAMP}_s \to \varnothing$ & $k_\mathrm{trk}[\mathrm{cAMP}_s]$ \\
11.E1 & $\varnothing \to \mathrm{PdiA}$ & $k_\mathrm{pdi}\,P(H)\,r(\mathrm{cAMP}_s)\ \text{(zeroth order)}$ \\
11.E2$^\ast$ & $\mathrm{PdiA} \to \varnothing$ & $k_\mathrm{deg}[\mathrm{PdiA}]$ \\
11.E3f$^\ast$ & $\mathrm{PdsA}_e+\mathrm{PdiA} \to \mathrm{PdsA{:}PdiA}$ & $k_{i,\mathrm{on}}[\mathrm{PdsA}_e][\mathrm{PdiA}]$ \\
11.E3r$^\ast$ & $\mathrm{PdsA{:}PdiA} \to \mathrm{PdsA}_e+\mathrm{PdiA}$ & $k_{i,\mathrm{off}}[\mathrm{PdsA{:}PdiA}]$ \\
11.E4$^\ast$ & $\mathrm{cAMP}_e+\mathrm{PdsA{:}PdiA} \to \mathrm{PdsA{:}PdiA}$ & $k_\mathrm{res}[\mathrm{cAMP}_e][\mathrm{PdsA{:}PdiA}]$ \\
11.2$^\ast$ & $\mathrm{cAMP}_e+\mathrm{PdsA}_e \to \mathrm{PdsA}_e$ & $k_\mathrm{pdsa}[\mathrm{cAMP}_e][\mathrm{PdsA}_e]$ \\
\bottomrule
\end{tabularx}
\end{table*}

\begin{table*}[!t]
\centering
\caption{Module~11 constants. $M=4817.6$ molecules\,$\mu\mathrm{M}^{-1}$ per voxel.}
\label{tab:m11const}
\footnotesize
\renewcommand{\arraystretch}{1.15}
\begin{tabularx}{\textwidth}{@{}l r l c >{\raggedright\arraybackslash}X@{}}
\toprule
Symbol & Value & Unit & Basis & Meaning, justification, source\\
\midrule
$k_\mathrm{pdsa}$ & \num{0.122} & $\mu\mathrm{M}^{-1}\mathrm{s}^{-1}$ & E & Effective second-order clearance constant, valid for cAMP well below $K_M$. \mkt{M11-pde-bader}{DdPDE1 (PdsA) is the main extracellular cAMP PDE}~\cite{bader2007}, with \mkt{M11-km}{an apparent $K_M$ of 0.8\,$\mu$M}~\cite{bader2007}. The value is set so that the clearance time and signal range of Sec.~\ref{sec:m11} result; it is not a measured $k_\mathrm{cat}/K_M$. \\
$k_\mathrm{bas},\ k_\mathrm{ind}$ & \num{4.48e-3},\ \num{4e-2} & $\mu\mathrm{M\,s^{-1}\,voxel^{-1}}$ & C & Constitutive and starvation-induced PdsA synthesis; $16\times$ the values that reproduce a $2\,\mu$M enzyme field at 1 cell per 36 voxels, so that the signal range $\lambda$ is $\approx\!1$--3 cell spacings and pulses are not global flashes (Sec.~\ref{sec:m11}). \\
$H_P,\ n_H$ & \num{25},\ 2 & $\mathrm{molec/voxel},\ -$ & E & Half-induction level and steepness of $P(H)$; the window opens early relative to the Hunger ramp. The \emph{pdsA} gene has \mkt{M11-promoters}{an aggregation-specific promoter}~\cite{faure1990} that is \mkt{M11-pa-off}{switched off after aggregation}~\cite{weening2003}; that later closing is outside the simulated time span. \\
$k_\mathrm{sec}$ & \num{0.01} & $\mathrm{s^{-1}}$ & E & Precursor-to-surface transfer ($\tau=100$\,s). \\
$k_\mathrm{shed}$ & \num{5e-4} & $\mathrm{s^{-1}}$ & E & Shedding of surface enzyme into the medium; equal to $k_\mathrm{deg}$, i.e.\ a 50:50 partition between shedding and loss. \\
$k_\mathrm{deg}$ & \num{5e-4} & $\mathrm{s^{-1}}$ & E & Turnover of all PdsA/PdiA protein states ($\tau=2000$\,s). This single constant sets the enzyme decay length $\sqrt{D/k_\mathrm{deg}}=361\,\mu$m. \\
$k_\mathrm{pdi}$ & \num{2.5e-3} & $\mu\mathrm{M\,s^{-1}\,voxel^{-1}}$ & E & Induced PdiA synthesis, co-induced with PdsA by the same starvation gate; \mkt{M11-inhib-absent}{PdiA is absent in growing cells and appears after transfer to starvation buffer}~\cite{franke1981}. \\
$K_\mathrm{pdi},\ n_\mathrm{pdi}$ & \num{0.05},\ 2 & $\mu\mathrm{M},\ -$ & E & Half-repression of PdiA synthesis by \mkt{M11-inhib-reg}{the slow cAMP signal cAMP$_s$}~\cite{yeh1978}. \\
$k_\mathrm{trk}$ & \num{4e-4} & $\mathrm{s^{-1}}$ & E & Unity-gain low-pass filter ($\tau=42$\,min) of cAMP$_e$ so that PdiA repression follows the pulse envelope. \\
$k_{i,\mathrm{on}},\ k_{i,\mathrm{off}}$ & \num{1},\ \num{1} & $\mu\mathrm{M}^{-1}\mathrm{s}^{-1},\ \mathrm{s^{-1}}$ & E & 1:1 binding with $K_d=1\,\mu$M. The measured stoichiometry is \mkt{M11-inhib-stoich}{1:1}~\cite{franke1981}, but the measured $K_d$ is \mkt{M11-kd}{$\approx\!10^{-10}$\,M}~\cite{franke1981}; the model's $K_d$ is $10^{4}$ times weaker to keep the equilibration fast against clearance. \\
$k_\mathrm{res}$ & \num{1.22e-4} & $\mu\mathrm{M}^{-1}\mathrm{s}^{-1}$ & E & Residual activity of the complex, $k_\mathrm{pdsa}/1000$. \\
$D_{\mathrm{cAMP}_e}$ & \num{350} & $\mu\mathrm{m^2/s}$ & D & \mkt{M11-dcamp}{Free-solution value $444$\,$\mu\mathrm{m^2/s}$ ($4.44\times10^{-6}\,\mathrm{cm^2/s}$)}~\cite{dworkin1977} reduced by a tortuosity factor; cells are not excluded volumes on the lattice. \\
$D_{\mathrm{PdsA}_e},D_{\mathrm{PdsA{:}PdiA}}$ & \num{65} & $\mu\mathrm{m^2/s}$ & E & Secreted enzyme (Diffusion set A of the model description). \\
$D_{\mathrm{PdiA}}$ & \num{60} & $\mu\mathrm{m^2/s}$ & E & Soluble inhibitor. \\
$D_{\mathrm{PdsA}_i},D_{\mathrm{PdsA}_m},D_{\mathrm{cAMP}_s}$ & \num{10},\ \num{0.1},\ \num{0} & $\mu\mathrm{m^2/s}$ & E & Precursor (cytosolic), surface form (membrane) and the local filter state (non-diffusing). \\
\bottomrule
\end{tabularx}
\end{table*}


\paragraph{Reactions.}
\begin{description}\itemsep0pt
\item[11.D1a/b] PdsA is synthesised in the cell at a constitutive rate (11.D1a; \mk{M11-pdsa-growth}{pdsA is already expressed during vegetative growth}~\cite{sucgang1997}) and at a starvation-induced rate gated by $P(H)$ (11.D1b), so that the enzyme appears with development. The gate depends only on the developmental variable $H$; the model contains no dependence of PdsA synthesis on pulses of cAMP.
\item[11.D2/D3] The enzyme is transferred from an intracellular precursor to the cell surface (11.D2) and shed into the medium (11.D3) as PdsA$_e$, which diffuses. \mk{M11-pdsa-er}{The intracellular pool of PdsA lies in the endoplasmic reticulum, a possible compartment for storage and secretion}~\cite{garcia2009}, and \mk{M11-pdsa-loc}{the enzyme can be secreted into the medium or exposed on the cell surface}~\cite{bader2007}. Only the secreted form clears extracellular cAMP; the surface-bound form is an explicit, non-diffusing species.
\item[11.D4] All protein states turn over with one rate constant; no measurement of the turnover of PdsA or PdiA is available, so 11.D4a--d and 11.E2 are design choices. The reactions for the secreted enzyme and the complex run in cells and on free voxels, so that clearance takes place everywhere.
\item[11.F1/F2] A per-voxel low-pass filter of cAMP$_e$, cAMP$_s$, lets PdiA repression follow the envelope of the pulse train and not individual pulses. Its time constant of 42\,min is a modelling choice; what it implements is \mk{M11-inhib-reg}{the repression of inhibitor synthesis by pulses of exogenous cAMP}~\cite{yeh1978}.
\item[11.E1/E2] PdiA is synthesised under the same starvation gate as PdsA and \mk{M11-inhib-reg}{repressed by cAMP$_s$}~\cite{yeh1978}, and degraded like PdsA.
\item[11.E3] PdsA and PdiA bind reversibly with \mk{M11-inhib-stoich}{1:1 stoichiometry}~\cite{franke1981}. The inhibition is titratable: as cells pack, the total pools scale with density, and the fraction of enzyme that is complexed rises.
\item[11.E4] The complex retains a small residual activity, $10^{-3}$ of the free enzyme (the factor is a choice); \mk{M11-complex-km}{the inhibitor raises the $K_M$ of the phosphodiesterase from about 10\,$\mu$M to 2\,mM}~\cite{franke1981}, so that the complex is not fully inactive.
\item[11.2] Free PdsA hydrolyses cAMP$_e$ (the enzyme is a catalyst and is not consumed); this is the primary clearance of the extracellular signal.
\end{description}

\paragraph{Constants and the signal range.}
Because PdsA$_e$ is produced per cell and lost per voxel, its steady-state level, and hence the clearance strength $k_\mathrm{eff}=k_\mathrm{pdsa}[\mathrm{PdsA}_e]$, is proportional to the cell density $\rho$.
The range over which a pulse signals, $\lambda=\sqrt{D_\mathrm{cAMP}/k_\mathrm{eff}}\propto\rho^{-1/2}$, scales like the mean cell spacing, so that $\lambda$ in units of cell spacings is independent of density for fixed rates, and the inhibitor titration breaks this only in the harmful direction: the fraction of complexed enzyme rises from 3\,\% at $\rho=7.7\times10^{-4}$ to 11\,\% at $6\times10^{-2}$ cells per voxel (it cannot exceed the ratio 0.12 of inhibitor to enzyme synthesis), and $\lambda$ in spacings grows as cells aggregate.
The synthesis constants $k_\mathrm{bas}$ and $k_\mathrm{ind}$ were raised by a factor of 16 above the values that reproduce a 2\,$\mu$M enzyme field at one cell per 36 voxels.
With the smaller values the clearance constant was $k_\mathrm{eff}=9.6\times10^{-4}\,\mathrm{s^{-1}}$, so that $\lambda\approx600\,\mu$m covered the whole population and the population fired as one synchronous global flash; at the present values $k_\mathrm{eff}=1.5\times10^{-2}\,\mathrm{s^{-1}}$ (clearance time 65\,s) and $\lambda\approx150\,\mu$m for a mean spacing of 50\,$\mu$m, i.e.\ about 1--3 cell spacings depending on the domain.
The enzyme field is initialised with its steady-state profile and not with a uniform value: the world boundary is absorbing, and $\sqrt{D_{\mathrm{PdsA}}/k_\mathrm{deg}}=361\,\mu$m is comparable to the domain, so that the boundary and not the degradation is the dominant sink, and the sustained field is 4.4--4.7 times higher at the centre than at the edge of the cell region.
The steady state depends on the geometry, so $\lambda$ has to be quoted per geometry.

Two constants deviate from measurements and should be read as model choices.
The measured dissociation constant of the PdsA--PdiA complex is \mk{M11-kd}{about $10^{-10}$\,M}~\cite{franke1981}; the model uses $1\,\mu$M, so that the binding equilibrates within seconds and clearance is not delayed.
And the effective clearance constant, $1.2\times10^{5}\,\mathrm{M^{-1}s^{-1}}$, is set by the required clearance time and is not a measured $k_\mathrm{cat}/K_M$; \mk{M11-km}{the apparent $K_M$ of DdPDE1 is $0.8\,\mu$M}~\cite{bader2007}, so the linear rate law overestimates clearance during pulses whose peak exceeds about $1\,\mu$M; the model compensates with a high enzyme concentration, and the surface pool of PdsA is $2.5\times10^{6}$ molecules per cell, which is a consequence of the slow shedding rate and the required flux and not a measured abundance.
The extracellular diffusion constant of cAMP is reduced from \mk{M11-dcamp}{the free-solution value of $4.44\times10^{-6}$\,cm$^2$/s}~\cite{dworkin1977} to $350\,\mu\mathrm{m^2/s}$, because cells are not excluded volumes on the lattice.

\subsection{Optional Module: Directional Memory}\label{sec:mmem}
This module is off in the default network and is described for completeness.
\mk{Mmem-mem-dir}{Cells keep their direction of motion in the back of a wave for the natural period of 6\,min, but increasingly reverse for longer periods}; \mk{Mmem-mem-carrier}{the molecular carrier of this memory is not known}~\cite{skoge2014}.
The optional module represents it phenomenologically by a membrane mark that is written where PIP3 is high (M.1) and erased slowly (M.2), and that does not feed back onto any other module; the motility reads it, so that cells steer on where PIP3 has been over the last minutes and not on its instantaneous value.
Because the mark is the time integral of the locally gated PIP3 with leakage $\tau=1/k_\mathrm{e}$, its centroid averages the PIP3 direction over the writing window and holds it after PIP3 falls.
No verified molecule is both PIP3-driven and long-lived, so the pool, the writing rate and the diffusion constant are estimates (Table~\ref{tab:mmemconst}); the erasure time of 3\,min is chosen to bridge the back of a wave.

\begin{table*}[!t]
\centering
\caption{Optional memory module (Mem), off by default and not used for the results of this paper: a phenomenological mark written where PIP3 is high and erased slowly. A coloured reaction is the one to which the equally coloured statement in the text and the equally coloured passage in the cited paper refer; a small square appends a further statement. $H(\mathrm{PIP3})=N_\mathrm{PIP3}^3/(K_\mathrm{mem}^3+N_\mathrm{PIP3}^3)$.}
\label{tab:mmemrxn}
\footnotesize
\renewcommand{\arraystretch}{1.15}
\begin{tabularx}{\textwidth}{@{}l>{\raggedright\arraybackslash}p{0.37\textwidth}>{\raggedright\arraybackslash}X@{}}
\toprule
ID & Reaction & Propensity (rate law)\\
\midrule
\mkt{Mmem-mem-carrier}{M.1} & \mkt{Mmem-mem-carrier}{$\mathrm{Mem}_c \to \mathrm{Mem}_m^*$} & $k_\mathrm{w}\,[\mathrm{Mem}_c]\,H(\mathrm{PIP3})$ \\
\mkn{Mmem-mem-dir}{Mmem-mem-carrier}{M.2} & \mkn{Mmem-mem-dir}{Mmem-mem-carrier}{$\mathrm{Mem}_m^* \to \mathrm{Mem}_c$} & $k_\mathrm{e}[\mathrm{Mem}_m^*]$ \\
\bottomrule
\end{tabularx}
\end{table*}

\begin{table*}[!t]
\centering
\caption{Constants of the optional memory module.}
\label{tab:mmemconst}
\footnotesize
\renewcommand{\arraystretch}{1.15}
\begin{tabularx}{\textwidth}{@{}l r l c >{\raggedright\arraybackslash}X@{}}
\toprule
Symbol & Value & Unit & Basis & Meaning, justification, source\\
\midrule
$[\mathrm{Mem}]_\mathrm{tot}$ & \num{5000} & $\mathrm{molec/voxel}$ & E & Large pool so that the centroid of the mark is not shot-noise limited. \\
$k_\mathrm{w}$ & \num{4e-3} & $\mathrm{s^{-1}}$ & E & Writing rate, chosen so that less than $\approx\!40\,\%$ of the pool is written per wave (an unsaturated integrator). \\
$K_\mathrm{mem},\ n_\mathrm{mem}$ & \num{5000},\ 3 & $\mathrm{molec/voxel},\ -$ & C & PIP3 level of half-writing, above resting PIP3 and inside the range reached on the rising edge of a wave. \\
$k_\mathrm{e}$ & \num{5.56e-3} & $\mathrm{s^{-1}}$ & E & Erasure with $\tau=3$\,min, so that the mark persists through the back of a wave: cells keep their direction there for waves of the natural 6\,min period~\cite{skoge2014}. \\
$D_{\mathrm{Mem}_c},D_{\mathrm{Mem}_m}$ & \num{10},\ \num{0.01} & $\mu\mathrm{m^2/s}$ & E & Cytosolic carrier and nearly immobile mark. \\
\bottomrule
\end{tabularx}
\end{table*}


